\documentclass[11pt]{article}
\usepackage[a4paper,margin=28mm]{geometry}
\usepackage{amsmath,amssymb,amsthm,booktabs}
\usepackage[hidelinks]{hyperref}
\usepackage[T1]{fontenc}
\usepackage{microtype}
\theoremstyle{definition}
\newtheorem{lemma}{Lemma}[section]
\newtheorem{proposition}[lemma]{Proposition}
\newtheorem{corollary}[lemma]{Corollary}
\newtheorem{definition}[lemma]{Definition}
\theoremstyle{remark}
\newtheorem{remark}[lemma]{Remark}
\theoremstyle{definition}
\newtheorem{theorem}[lemma]{Theorem}
\numberwithin{equation}{section}
\newcommand{\Q}{\mathbb Q}\newcommand{\C}{\mathbb C}
\newcommand{\R}{\mathbb R}\newcommand{\Z}{\mathbb Z}\newcommand{\N}{\mathbb N}
\newcommand{\PP}{\mathbb P}
\newcommand{\eps}{\varepsilon}\newcommand{\card}{\#}
\newcommand{\ST}{\mathsf{ST}}\newcommand{\CC}{\mathsf{CC}}\newcommand{\GE}{\mathsf{GE}}
\newcommand{\lean}{Lean}
\newcommand{\OO}{\mathcal O}
\newcommand{\Nm}{\operatorname{N}}
\newcommand{\hgt}{\operatorname{ht}}
\renewcommand{\labelitemi}{--}
\setlength{\emergencystretch}{3em}
\title{Powerful values of quadratic polynomials}
\author{Danny Plymesser}
\date{16 November 2026\\[0.3em]{\small
\href{https://doi.org/10.5281/zenodo.24292371}{doi: 10.5281/zenodo.24292371}}}
\hypersetup{pdftitle={Powerful values of quadratic polynomials},
pdfauthor={Danny Plymesser},pdfsubject={DOI: 10.5281/zenodo.24292371}}
\begin{document}\maketitle\vspace{-1.5em}
\begin{abstract}
Let $f$ be a quadratic polynomial with integer coefficients, positive
leading coefficient and nonsquare discriminant. We show that the number of
$n\le N$ with $f(n)$ powerful is $O(N^{\vartheta})$ for every
$\vartheta>293/1220$. A second bound has the exponent $389/1600$ and a constant
that grows only polynomially with the coefficients of $f$; it is applied to
pairs consisting of a $6$-full or $8$-full integer and a neighbouring
powerful integer. A fixed-quadratic substitution also gives a bound for
$1+a^3t^3$. The proofs factor $f(n)$ in the quadratic
field of its discriminant, apply a determinant method with two contact
polynomials to the resulting equation of bidegree three and two, and count
the points on the auxiliary curves with the lifted-curve estimate
of~\cite{powerful}. The main quadratic and mixed-full bounds are also
formalised in \lean, with the three external estimates recorded explicitly.
\end{abstract}

\section{Introduction}\label{sec:intro}
A positive integer $m$ is \emph{powerful} if $p^2\mid m$ for every prime
$p\mid m$, and \emph{$k$-full} if $p^k\mid m$ for every prime $p\mid m$.
There are about $x^{1/2}$ powerful integers up to $x$. For integers $a,b,c$
put
\[
 S(a,b,c;N)=\card\{1\le n\le N:\ an^2+bn+c\ge1,\ an^2+bn+c\textnormal{ powerful}\}.
\]
The values $an^2+bn+c$ with $n\le N$ have size $N^2$, so counting powerful
integers gives only $S(a,b,c;N)\ll N$. The elementary count of
Section~\ref{sec:easy} gives $N^{2/5+\eps}$, and the true order is expected
to be far smaller. This paper proves upper bounds for $S(a,b,c;N)$ when the
discriminant $b^2-4ac$ is not a square. It develops the counting method
of~\cite{powerful} for polynomial values, including families in which the
coefficients vary.

\subsection{Quadratic estimates}\label{sec:results}
\begin{theorem}[fixed quadratic]\label{thm:fixed}
Let $A,B,C\in\Z$ with $A>0$ and $B^2-4AC$ not a square.
For every $\eps>0$ there is $K$ such that
\[
 S(A,B,C;N)\le KN^{293/1220+\eps}\qquad(N\ge1).
\]
\end{theorem}
Here $K$ depends on $A,B,C,\eps$, and $N$ is real. Negative discriminants
are allowed.

\begin{theorem}[polynomial dependence on the height]\label{thm:uniform}
There are absolute constants $K,\nu>0$ such that for
all $a,b,c\in\Z$ with $a\ge1$ and $b^2-4ac$ not a square, every real
$L\ge2$ with $|a|,|b|,|c|\le L$, and every $N\ge2$,
\[
 S(a,b,c;N)\le KL^{\nu}N^{389/1600}.
\]
Moreover there are absolute $\eta,C>0$ such that, with
$M_0=\lfloor N^{\eta}\rfloor$,
\[
 \sum_{a=1}^{M_0}\ \sum_{|b|\le M_0}\ \sum_{\substack{|c|\le M_0\\ b^2-4ac\textnormal{ not a square}}}
 S(a,b,c;N)\le CN^{390/1600}\qquad(N\ge2).
\]
\end{theorem}
Neither bound has an $\eps$. The exponent $\nu$ is not computed. The second
bound follows by summing the first; one may take
$\eta=1/(1600(\nu+3))$. See
Sections~\ref{sec:A3proof} and~\ref{sec:A3remarks}.

\begin{theorem}[the family $n^2+a^2$]\label{thm:family}
Put
\[
 \Theta(\kappa)=\begin{cases}
 (293+153\kappa)/1220,&\kappa\le1/28,\\
 (81+61\kappa)/340,&1/28<\kappa\le19/73,\\
 (61+105\kappa)/310,&\kappa>19/73.
 \end{cases}
\]
For every $0\le\kappa\le289/1050$ and $\eps>0$ there is $K$ such that, for
all $N\ge2$ and all $a\in\Z$ with $1\le|a|\le N^{\kappa}$,
\[
 \card\{1\le n\le N:n^2+a^2\textnormal{ is powerful}\}\le KN^{\Theta(\kappa)+\eps}.
\]
\end{theorem}
The function $\Theta$ is continuous, with $\Theta(0)=293/1220$ and
$\Theta(289/1050)=29/100$.

\subsection{Cubic values and full pairs}
\begin{corollary}[cubic values]\label{cor:cubic}
There is an absolute $\nu_1\ge0$ such that for every $\eps>0$
there is $C$ with
\[
 \card\{1\le t\le T:1+At^3\textnormal{ is powerful}\}\le CA^{\nu_1}T^{3/5+\eps}
 \qquad(A\ge1,\ T\ge2).
\]
For every $\eps>0$ there is also $C_\eps$ such that
\[
 \card\{1\le t\le T:1+a^3t^3\textnormal{ is powerful}\}
 \le C_\eps(aT)^{293/1220+\eps}\qquad(a\ge1,\ T\ge2).
\]
\end{corollary}

Here $A,a$ are positive integers and $T$ is real.

In the following two theorems $s\in\{6,8\}$, and for real $z>0$
\[
 \delta_s(z)=\min\Bigl\{\frac{1}{6400(z+1)s(s+1)},\ \frac12\Bigl(\frac1s-\frac{49}{400}\Bigr)\Bigr\},
\]
which is positive. The trivial bound for the number of $s$-full $n\le Y$ is
$\ll Y^{1/s}$.

\begin{theorem}[a positive shift]\label{thm:positive}
For every integer $h\ge1$ there are $K_h,C_h\ge1$ such
that
\[
 \card\{1\le t\le T:At^2+h\textnormal{ is powerful}\}\le K_hA^{C_h}T^{389/1600}
 \qquad(A\ge1,\ T\ge2),
\]
and for all $s\in\{6,8\}$, $2\le r<s$ and $\eps>0$ there is $C$ with
\[
 \card\{1\le n\le Y:n\textnormal{ is }s\textnormal{-full},\ n+h\textnormal{ is }r\textnormal{-full}\}
 \le CY^{1/s-\delta_s(C_h)+\eps}\qquad(Y\ge2).
\]
\end{theorem}

\begin{theorem}[the reverse neighbour]\label{thm:reverse}
There are functions $K_*,\nu_*:\N\to\R$, with
$K_*(H')>0$ and $\nu_*(H')>0$ for every $H'\ge1$, such that for all
$A,H'\ge1$ with $AH'$ not a square and all $T\ge2$,
\[
 \card\{1\le t\le T:At^2-H'\ge1,\ At^2-H'\textnormal{ is powerful}\}
 \le K_*(H')A^{\nu_*(H')}T^{389/1600},
\]
and for all $s\in\{6,8\}$, $2\le r<s$ and $\eps>0$ there is $C$ with
\[
 \card\{2\le n\le Y:n\textnormal{ is }s\textnormal{-full},\ n-1\textnormal{ is }r\textnormal{-full}\}
 \le CY^{1/s-\delta_s(\nu_*(1))+\eps}\qquad(Y\ge2).
\]
\end{theorem}
The proofs use only that the neighbour is powerful
(Remark~\ref{rem:rs}).

The coefficient exponents in these applications can be chosen from a
single constant: Section~\ref{sec:consequences} takes
$C_h=\max(1,\nu)$ and $\nu_*(H')=\nu$, with $\nu$ from
Theorem~\ref{thm:uniform}. Thus the shifts affect $K_h$ and $K_*(H')$,
but need not affect these exponents. The separate cubic argument in
Section~\ref{sec:A4} permits the absolute choice $\nu_1=180$.

\subsection{Comparison with earlier work}\label{sec:compare}
For a fixed quadratic with nonzero discriminant,
Wongcharoenbhorn--Meemark~\cite[Theorem~B]{wm} prove the exponent
$\approx0.4769$, and the unrefereed manuscript~\cite[Theorem~1.3]{saito}
states $2/5$. Both allow square discriminants, which
Theorem~\ref{thm:fixed} excludes. For $n^2+a^2$,
\cite[Theorem~A]{wm} gives $29/100$ uniformly for
$1\le|a|\le N^{29/100}$; for $n^2+1$, Puchta~\cite{puchta} proved a bound
of order $N^{2/5}$ up to logarithmic factors.
Theorem~\ref{thm:fixed} gives $293/1220$ for fixed $a\ne0$.
Theorem~\ref{thm:family} improves $29/100$ in the shorter range
$1\le|a|\le N^\kappa$, $\kappa<289/1050$, with equality at the endpoint.

Polynomial coefficient dependence alone is not new.
Evertse--Silverman~\cite[Theorem~1]{es} and
Vaughan--Wooley~\cite[Lemma~3.5]{vw}, with elementary class-number bounds,
give $S(a,b,c;N)\ll_\eps L^{39}N^{2/5+\eps}$ in the range of
Theorem~\ref{thm:uniform}. Appendix~\ref{app:classical} derives this coarse
coefficient power. The distinction in Theorem~\ref{thm:uniform} is the
exponent $389/1600$ with an absolute polynomial height loss; we found no
earlier bound combining these features for all the stated quadratics.
The coefficient average in~\cite[Remark~4.5]{wm}, over heights at least
$N^{2+\eps}$, concerns a different range.

The cubic bounds are applications. The general estimate follows from
Evertse--Silverman; the cube-coefficient estimate follows from
Theorem~\ref{thm:fixed} by factoring $u^3+1$. The same reduction
and~\cite{saito} give $(aT)^{2/5+\eps}$. For the mixed-full applications,
the cases $(r,s)=(2,6),(3,6),(2,8),(3,8),(4,8)$ improve $Y^{1/s}$;
we found no stronger literature bound for these cases with the stated
shifts and orientation. The other pairs illustrate the method without
a best-known-bound claim.

Related results address different questions: infinitely many powerful
$n^2+k$ for almost all $k$~\cite{ddl}, long blocks under Vojta's
conjecture~\cite{pasten}, and counts of norm-form or $S$-squareful values
in several variables~\cite{streeter,pohl}. The latter do not yield our
one-variable estimates by fixing a coordinate. Under $abc$ the quadratic
count is $\ll N^\eps$~\cite[Proposition~3.7]{wm}; no such conjecture is
assumed here. These comparisons are not an exhaustive literature search.

\subsection{Method and scope}\label{sec:notclaimed}
Put $\Delta=B^2-4AC$. The element $(2An+B)+\sqrt\Delta$ has norm
$4Af(n)$. Factoring its ideal, choosing class representatives and
balancing conjugates reduces a powerful value to
$\gamma\mathsf Y^3\mathsf X^2=(2An+B)+\sqrt\Delta$.
In primitive coordinates $y,x\in\Z^2$ this gives
\begin{equation}\label{eq:introF}
 F(y;x)=e\ne0,\qquad \deg F=(3,2).
\end{equation}
This replaces the equation $e^3v^2-d^3u^2=h$ in~\cite{powerful}.
Thue and Pell estimates handle boxes with a short side. In the other
boxes the slopes lie within $O(1/n)$ of a real curve.
Heath-Brown's determinant method~\cite{hbannals,hb} places the points in
each slope interval of length $X^{-\mu}$ on an auxiliary curve when
$\mu>57/305$. The threshold comes from two contact polynomials of
bidegrees $(2,1)$ and $(1,2)$ and fifth-order vanishing.
Low-degree components have $X^{o(1)}$ points; the remaining components
are counted by the lifted-curve lemma of~\cite{powerful}, using the
affine-curve bound of~\cite{ccdn}. Lattice moments then give
$X^{1/10+3\mu/4+o(1)}$, with limiting exponent $293/1220$.

Sections~\ref{sec:charts}--\ref{sec:sum} carry this out for a fixed
quadratic. Section~\ref{sec:uniform} makes the coefficient losses
polynomial, using bounded-norm balancing elements and quantitative graph
charts. Section~\ref{sec:A2} exploits the Gaussian structure of $n^2+a^2$.
The applications follow in Sections~\ref{sec:A4} and~\ref{sec:consequences}.
Split quadratics such as $n(n+1)$ are excluded here and treated
in~\cite{powerful}. The fixed-quadratic constant is not effective; the
uniform coefficient exponent is absolute but large and uncomputed.
The coefficient average and mixed-full savings are therefore asymptotic,
not practical numerical estimates. The exponent $389/1600$ is not optimized.

The main quadratic and mixed-full bounds also have proofs in \lean~4
with Mathlib. Three published estimates, stated in
Section~\ref{sec:cited}, enter that formalisation as explicit hypotheses.
Their proofs lie outside the \lean\ development, so its conclusions are
conditional on these inputs. The mathematical arguments invoke the
published theorems. Section~\ref{sec:status}
and Appendix~\ref{app:lean} describe the verification and authorship and
map the paper's statements to the formal interfaces.

\section{Notation and results from the literature}\label{sec:cited}
Counted variables are positive integers unless stated otherwise. An
unspecified constant may depend on the parameters that precede it in the
statement and on nothing else. We write $\omega(m)$ for the number of
distinct prime divisors of $m$, $\tau(m)$ for the number of divisors, and
$\|G\|$ for the largest absolute value of a coefficient of a polynomial
$G$. A vector $y=(y_0,y_1)\in\Z^2$ is \emph{primitive} if
$\gcd(y_0,y_1)=1$, and $\hgt(y)=\max(|y_0|,|y_1|)$. A polynomial in two
pairs of variables is \emph{bihomogeneous of bidegree $(\alpha,\beta)$} if
it is homogeneous of degree $\alpha$ in the first pair and of degree
$\beta$ in the second.

We use the following three published estimates throughout. They are stated
in the forms needed here. In the \lean\ development these statements are
explicit hypotheses; the external theorems themselves are not formalised
in this repository.

\paragraph{$\ST$ (Stewart \cite[Theorem~1 and p.~796]{stewart}).}
Let $F\in\Z[u,v]$ be a binary form of degree $\deg F\ge3$ with coprime
coefficients that is squarefree in $\Q[u,v]$. Then for every nonzero
$m\in\Z$,
\[
 \card\{(u,v)\in\Z^2:\gcd(u,v)=1,\ F(u,v)=m\}\le2800\,(\deg F)^{1+\omega(|m|)}.
\]

\paragraph{$\CC$ (Castryck--Cluckers--Dittmann--Nguyen \cite[Theorem~3]{ccdn}).}
For every integer $w\ge2$ there is $c_w$ with the following property. Let
$I\subset\Q[z_1,\dots,z_w]$ be a prime ideal defining an affine curve of
degree $d\ge1$; that is, the quotient $\Q[z_1,\dots,z_w]/I$ has Krull
dimension one and, for all large $k$, the polynomials of degree at most $k$
span a subspace of dimension $dk+k_0$ in it, for some integer $k_0$. Then for
every $\Omega\ge1$ the number of integer zeros $P$ of $I$ with
$\max_i|P_i|\le\Omega$ is at most
\[
 c_w\,d^3\Omega^{1/d}(\log\Omega+d).
\]

\paragraph{$\GE$ (Gelfond; \cite[Lemma~1.6.11]{bg}).}
For every $w$ there is $c\ge0$ such that if $F,G\in\Z[z_1,\dots,z_w]$,
$F\ne0$ and $G\mid F$, then $\|G\|\le\exp(c\deg F)\,\|F\|$, where $\deg F$
is the total degree.

\medskip
In $\ST$ reducible forms are allowed, and $(u,v)$ and $(-u,-v)$ are counted
separately. In $\CC$ the curve need not be geometrically irreducible. We
use $\CC$ only for $w=12$ and $\GE$ only for $w=4$.
Section~\ref{sec:status} describes the comparison with the sources.

\section{Arithmetic charts}\label{sec:charts}
Throughout Sections~\ref{sec:charts}--\ref{sec:sum}, fix integers
$A>0,B,C$ with nonsquare $\Delta=B^2-4AC$, and let $f(n)=An^2+Bn+C$;
constants may depend on them. Write $\theta=\sqrt\Delta$ and
$\Nm(a+b\theta)=a^2-\Delta b^2$, which vanishes only at zero.
We call $a,b$ the rational and $\sqrt\Delta$-coordinates.
Using the finite class group and, for $\Delta>0$, a norm-one unit,
we construct finitely many integral forms $F_\chi$ of bidegree $(3,2)$
whose primitive solutions $F_\chi(y;x)=e_\chi$ cover the powerful
values. No hypothesis from the literature is used here.

\subsection{Factorisation in the maximal order}
Let $\OO$ be the maximal order of $\Q(\sqrt\Delta)$. For $n\ge1$ put
$\alpha_n=(2An+B)+\theta$, which lies in $\OO$ because
$\Delta\equiv B^2\pmod4$, and
$\bar\alpha_n=(2An+B)-\theta$. Then
\begin{equation}\label{eq:alphanorm}
 \alpha_n\bar\alpha_n=(2An+B)^2-\Delta=4A\,f(n).
\end{equation}
Every powerful integer is $d_0^3u_0^2$ with $d_0,u_0\ge1$: take
$d_0$ to be the product of primes occurring to odd exponent.

\begin{lemma}\label{lem:ideal}
Let $n\ge1$ with $f(n)>0$ and $f(n)$ powerful. Then
$(\alpha_n)=R\,\mathcal A^3\mathcal B^2$ with nonzero ideals
$R,\mathcal A,\mathcal B$ of $\OO$ such that $R$ divides $(4A\Delta)$.
\end{lemma}
\begin{proof}
Write $f(n)=d_0^3u_0^2$. If $v_P(\alpha_n)=1$, then either
$\bar\alpha_n\in P$, giving $2\theta\in P$ and $P\mid(4\Delta)$, or
$v_P(\bar\alpha_n)=0$. In the latter case $P\nmid(4A)$ would give
$1=3v_P(d_0)+2v_P(u_0)$ by \eqref{eq:alphanorm}, a contradiction.
Thus all exponent-one prime ideals divide $(4A\Delta)$. Put these into
$R$, allocate an odd exponent $e\ge3$ as $3+2((e-3)/2)$ to
$\mathcal A^3\mathcal B^2$, and allocate even exponents to $\mathcal B^2$.
\end{proof}

Fix integral ideal representatives $\mathcal R$ of the finite class
group and nonzero $\rho_I\in I$ for each $I\in\mathcal R$.
Let $\mathcal J$ consist of triples $\pi=(R,I,J)$ with
$R\mid(4A\Delta)$ and $I,J\in\mathcal R$. Choose a generator
$g_\pi$ of $RI^3J^2$ when principal, and put $g_\pi=1$ otherwise;
set $c_\pi=g_\pi/(\rho_I^3\rho_J^2)\ne0$.

\begin{lemma}\label{lem:unbalanced}
For every $n\ge1$ with $f(n)>0$ powerful there are $\pi\in\mathcal J$ and
nonzero $\mathsf Y,\mathsf X\in\OO$ with $\alpha_n=c_\pi \mathsf Y^3\mathsf X^2$.
\end{lemma}
\begin{proof}
Write $\mathcal A=(\mathsf Y_0)I$, $\mathcal B=(\mathsf X_0)J$ using
Lemma~\ref{lem:ideal} and the class representatives. Then $RI^3J^2$ is
principal and $\alpha_n=\epsilon g_\pi\mathsf Y_0^3\mathsf X_0^2$ for a
unit $\epsilon$. Set $\mathsf Y=\epsilon\mathsf Y_0\rho_I$ and
$\mathsf X=\epsilon^{-1}\mathsf X_0\rho_J$. These are integral since
$\mathsf Y_0\rho_I\in\mathcal A$, $\mathsf X_0\rho_J\in\mathcal B$;
the powers of $\epsilon$ give the required identity.
\end{proof}

Fix $m_{\OO}\ge1$ with $\OO\subset m_{\OO}^{-1}\Z[\theta]$ and put
$\iota(y)=(y_0+y_1\theta)/m_{\OO}$ for $y\in\Z^2$.
Then $m_{\OO}^2\Nm\iota(y)=y_0^2-\Delta y_1^2$ is integral,
and nonzero for $y\ne0$.

\subsection{Balancing}
Put $\|a+b\theta\|^2=a^2+|\Delta|b^2$. For $\Delta<0$ this is
$\Nm z$; for $\Delta>0$ we normalize by units. For $\Delta>0$, direct expansion gives
\begin{equation}\label{eq:parallelogram}
 2\|z\|^2\|w\|^2=\|zw\|^2+\|\bar zw\|^2 .
\end{equation}
Consequently $\|zw\|^2\le2\|z\|^2\|w\|^2$,
$\|z^3\|^2\le4\|z\|^6$, and, using $\bar z(zw)=(\Nm z)w$,
$(\Nm z)^2\|w\|^2\le2\|z\|^2\|zw\|^2$.
Also $\|z^2\|^2-\|z\|^4=4\Delta a^2b^2\ge0$.

\begin{lemma}\label{lem:balance}
Let $\Delta>0$. There are a unit $\epsilon_0$ of $\OO$ of norm $1$ and
$U_0\ge1$ such that for every nonzero $\mathsf Y\in\OO$ there is $k\in\Z$ with
$\|\epsilon_0^{2k}\mathsf Y\|^2\le U_0\,|\Nm \mathsf Y|$.
\end{lemma}
\begin{proof}
Choose a Pell unit $\epsilon_0=x_1+y_1\theta$ with integers $x_1>1$,
$y_1>0$ and norm $1$. Write $\epsilon_0^2=U_0+V_0\theta$, where
$U_0>1$, $V_0>0$ and $U_0^2-\Delta V_0^2=1$. Minimize the nonnegative
integer $m_{\OO}^2\|\epsilon_0^{2k}\mathsf Y\|^2$ over $k\in\Z$,
and write the minimizing element as $a+b\theta$, with
$s_+=a^2+\Delta b^2$ and $s_-=a^2-\Delta b^2=\Nm\mathsf Y$.
The neighboring values are
$(U_0^2+\Delta V_0^2)s_+\pm4U_0V_0\Delta ab$; minimality gives
$V_0s_+\pm2U_0ab\ge0$. Therefore
\[
 U_0^2s_-^2-s_+^2
 =\Delta(V_0s_+-2U_0ab)(V_0s_++2U_0ab)\ge0,
\]
which proves the bound.
\end{proof}

\begin{lemma}\label{lem:growth}
Let $n_0=|B|+2|\Delta|+1$ and $G_0=(2A+|B|)^2+|\Delta|+1$. For $n\ge n_0$,
\[
 n^2\le2\,\Nm\alpha_n,\qquad \Nm\alpha_n\le G_0n^2,\qquad
 \|\alpha_n\|^2\le3\,\Nm\alpha_n ,
\]
and hence $f(n)>0$.
\end{lemma}
\begin{proof}
Let $m=2An+B$. Since $A\ge1$ and $n\ge|B|$ we have $n\le m\le(2A+|B|)n$, and
$m^2\ge n^2\ge2|\Delta|$. The three bounds follow from
$\Nm\alpha_n=m^2-\Delta$ and $\|\alpha_n\|^2=m^2+|\Delta|$.
\end{proof}

\begin{lemma}\label{lem:balancedcover}
There is $U_1\ge1$ such that for every $n\ge n_0$ with $f(n)$ powerful there
are $\pi\in\mathcal J$ and nonzero $\mathsf Y,\mathsf X\in\OO$ with $\alpha_n=c_\pi \mathsf Y^3\mathsf X^2$,
$\|\mathsf Y\|^2\le U_1|\Nm \mathsf Y|$ and $\|\mathsf X\|^2\le U_1|\Nm \mathsf X|$.
\end{lemma}
\begin{proof}
For $\Delta<0$ use Lemma~\ref{lem:unbalanced} and $U_1=1$.
For $\Delta>0$, replace its pair $(\mathsf Y,\mathsf X)$ by
$(\epsilon_0^{2k}\mathsf Y,\epsilon_0^{-3k}\mathsf X)$ using
Lemma~\ref{lem:balance}. The product is unchanged and
$\|\mathsf Y\|^2\le U_0|\Nm\mathsf Y|$. With $c=c_\pi$,
\eqref{eq:parallelogram} and Lemma~\ref{lem:growth} give
\[
 (\Nm(c\mathsf Y^3))^2\|\mathsf X\|^4
 \le2\|c\mathsf Y^3\|^2\|\alpha_n\|^2
 \le48\|c\|^2U_0^3|\Nm\mathsf Y|^3|\Nm\alpha_n|.
\]
Substitute $|\Nm\alpha_n|=|\Nm c|\,|\Nm\mathsf Y|^3|\Nm\mathsf X|^2$
and cancel to obtain
$|\Nm c|\,\|\mathsf X\|^4\le48\|c\|^2U_0^3|\Nm\mathsf X|^2$.
One may take $U_1=U_0+1+\sum_\pi48\|c_\pi\|^2U_0^3/|\Nm c_\pi|$.
\end{proof}

\subsection{The charts}
Clear denominators in the $\sqrt\Delta$-coordinate of
$c_\pi\iota(y)^3\iota(x)^2$ to obtain an integral form $F_\pi$
of bidegree $(3,2)$ and a denominator $e_\pi\ge1$.
A \emph{chart} is $\chi=(\pi,g,g')$ with $\pi\in\mathcal J$,
$1\le g,g'\le G_1=1+\sum_\pi e_\pi$. Set
$\gamma_\chi=g^3g'^2c_\pi$, $e_\chi=e_\pi$ and
$F_\chi=g^3g'^2F_\pi$, so
\begin{equation}\label{eq:coeffid}
 F_\chi(y;x)=e_\chi\cdot\bigl(\sqrt\Delta\textnormal{-coordinate of }
 \gamma_\chi\,\iota(y)^3\iota(x)^2\bigr).
\end{equation}
The set $\mathcal X$ of charts is finite. The form $F_\chi$ is not zero:
otherwise \eqref{eq:coeffid} at $\iota(y)=\iota(x)=1$ and at
$\iota(y)=\theta$, $\iota(x)=1$ (where $\theta^3=\Delta\theta$) would show
that both coordinates of $\gamma_\chi$ vanish.

\begin{definition}\label{def:point}
Let $H\ge1$, $\chi\in\mathcal X$ and $n\ge1$. A pair $(y,x)$ of primitive
vectors in $\Z^2$ is an \emph{arithmetic point of $\chi$ over $n$} if
\begin{equation}\label{eq:element}
 \gamma_\chi\,\iota(y)^3\iota(x)^2=(2An+B)+\sqrt\Delta
\end{equation}
and, with $D_y=|\Nm\iota(y)|$ and $D_x=|\Nm\iota(x)|$,
\[
 H^{-1}\le D_y,\quad H^{-1}D_y\le\hgt(y)^2\le HD_y,\quad
 H^{-1}\le D_x,\quad H^{-1}D_x\le\hgt(x)^2\le HD_x,
\]
\[
 H^{-1}n^2\le D_y^3D_x^2\le Hn^2 .
\]
\end{definition}
By \eqref{eq:coeffid} and \eqref{eq:element} an arithmetic point satisfies
$F_\chi(y;x)=e_\chi$.

\begin{proposition}\label{prop:charts}
There is $H\ge1$ such that for every $n\ge n_0$ with $f(n)$ powerful there
are a chart $\chi\in\mathcal X$ and an arithmetic point of $\chi$ over $n$.
\end{proposition}
\begin{proof}
Represent the balanced pair of Lemma~\ref{lem:balancedcover} by
$y',x'\in\Z^2$. Since $F_\pi(y';x')=e_\pi$, their contents $g,g'$ satisfy
$g^3g'^2\mid e_\pi$. Thus $\chi=(\pi,g,g')\in\mathcal X$, and the
primitive vectors $y=y'/g$, $x=x'/g'$ satisfy \eqref{eq:element}.
Balancing is unchanged by scaling. Moreover
\[
 1\le m_{\OO}^2D_y=|y_0^2-\Delta y_1^2|
 \le(1+|\Delta|)\hgt(y)^2,
 \qquad \hgt(y)^2\le m_{\OO}^2U_1D_y,
\]
and likewise for $x$. Taking norms gives
$|\Nm\gamma_\chi|D_y^3D_x^2\in[n^2/2,G_0n^2]$ by
Lemma~\ref{lem:growth}. A sufficiently large $H$, uniform over the
finite chart list, gives all the required inequalities.
\end{proof}

\begin{lemma}\label{lem:recover}
If $(y,x)$ is an arithmetic point of $\chi$ over $n$ and over $n'$, then
$n=n'$; and $\hgt(y),\hgt(x)\le H^3n$. Consequently, for $X\ge1$ let
$\mathcal W(X)$ be the set of $n\ge n_0$ with $X<n\le2X$ and $f(n)$
powerful, and $\mathcal S(X)$ the set of all triples $(\chi,y,x)$ such that
$(y,x)$ is an arithmetic point of $\chi$ over some $n\in\mathcal W(X)$. Then
$\mathcal S(X)$ is finite and $\card\mathcal W(X)\le\card\mathcal S(X)$.
\end{lemma}
\begin{proof}
The rational coordinate in \eqref{eq:element} determines $n$ since
$A\ne0$. As $HD_y,HD_x\ge1$,
$\hgt(y)^2\le HD_y\le H^5D_y^3D_x^2\le H^6n^2$, and likewise for $x$.
Thus $\mathcal S(X)$ is finite; choosing a triple for each $n$ by
Proposition~\ref{prop:charts} gives an injection.
\end{proof}
A given $n$ may be represented by several triples; all bounds below are
upper bounds for the number of triples.

\subsection{Slopes and the real curve}
Write $\sigma_\pm(a+b\theta)=a\pm b\sqrt\Delta$.
A \emph{dominant coordinate} $d_y$ of primitive $y$ has
$|d_y|=\hgt(y)$; the other coordinate divided by $d_y$ is its
\emph{slope} $\xi\in[-1,1]$. A \emph{frame} $\phi$ chooses a
chart $\chi$ and dominant indices for $y,x$, giving four frames per
chart. Then $\sigma_\pm(\iota(y))=d_y\ell_\pm(\xi)$ and
$\sigma_\pm(\iota(x))=d_xm_\pm(\zeta)$ for fixed linear forms
(e.g. $\ell_\pm(\xi)=(\xi\pm\sqrt\Delta)/m_{\OO}$ when $d_y=y_1$).
Put
\[
 R_\phi(\xi,\zeta)=\frac{\sigma_+(\gamma_\chi)\ell_+(\xi)^3m_+(\zeta)^2
 -\sigma_-(\gamma_\chi)\ell_-(\xi)^3m_-(\zeta)^2}{2\sqrt\Delta}.
\]
This real polynomial has degrees at most $(3,2)$ and
$F_\chi(y;x)=e_\chi d_y^3d_x^2R_\phi(\xi,\zeta)$.
Choose $M_1\ge1$ bounding all $|\ell_\pm|,|m_\pm|$ on $[-1,1]$
and set $\eta_0=(HM_1)^{-1}$.

\begin{lemma}\label{lem:slopes}
Let $(y,x)$ be an arithmetic point of $\chi$ over $n$, in a frame $\phi$ of
$\chi$ given by dominant coordinates $d_y,d_x$, with slopes $(\xi,\zeta)$.
Then $|\ell_\pm(\xi)|\ge\eta_0$, $|m_\pm(\zeta)|\ge\eta_0$,
\[
 R_\phi(\xi,\zeta)=\frac1{d_y^3d_x^2}\qquad\textnormal{and}\qquad
 |R_\phi(\xi,\zeta)|\le\frac{H^3}{n}.
\]
\end{lemma}
\begin{proof}
The product $|\ell_+(\xi)\ell_-(\xi)|=D_y/\hgt(y)^2\ge H^{-1}$
and the upper bound $M_1$ give the lower bound for each factor;
the same applies to $m_\pm$. Subtract the two embeddings of
\eqref{eq:element} and divide by $2\sqrt\Delta\,d_y^3d_x^2$ for the
identity. Finally
$n^2\le HD_y^3D_x^2\le H^6d_y^6d_x^4$.
\end{proof}
The set of $(\xi,\zeta)\in[-1,1]^2$ with
$|\ell_\pm(\xi)|,|m_\pm(\zeta)|\ge\eta_0$ is the \emph{balanced patch} of
$\phi$.

\subsection{Boxes and windows}
For a triple $(\chi,y,x)$ put $U=H|\Nm\iota(y)|$ and $V=H|\Nm\iota(x)|$.
\begin{lemma}\label{lem:norms}
If $(y,x)$ is an arithmetic point of $\chi$ over $n\le2X$, then $U\ge1$,
$V\ge1$ and $U^3V^2\le4H^6X^2$.
\end{lemma}
\begin{proof}
$U^3V^2=H^5D_y^3D_x^2\le H^6n^2\le4H^6X^2$.
\end{proof}

\begin{lemma}\label{lem:partition}
Let $A_0\ge1$. There is $K_{\mathrm p}>0$ such that for every $X\ge2$ and every finite
set $\mathcal S$ with functions $U,V:\mathcal S\to[1,\infty)$ satisfying
$U^3V^2\le A_0X^2$, the sets
\[
 \mathcal S(D,E)=\{z\in\mathcal S:\ D/2<U(z)\le D,\ E/2<V(z)\le E\},
\]
with $D,E$ running over the powers $2^j$, $j\ge0$, partition $\mathcal S$;
at most $K_{\mathrm p}(1+\log X)^2$ of them are nonempty; and each nonempty one has
$D^3E^2\le32A_0X^2$.
\end{lemma}
\begin{proof}
Each $z$ lies in exactly one set. If $z\in\mathcal S(D,E)$ then
$D<2U\le2U^3V^2\le2A_0X^2$ and likewise for $E$, so there are at most
$(2+\log_2(2A_0X^2))^2$ nonempty sets, and
$D^3E^2\le2^3U^3\cdot2^2V^2\le32A_0X^2$.
\end{proof}
We use this with $A_0=4H^6$, so that $D^3E^2\le128H^6X^2$; the pair $(D,E)$
is called a \emph{box}.

\begin{lemma}\label{lem:prefix}
Let $\vartheta>0$, $C_{\mathrm w}\ge0$, $X_0\ge1$, $n_1\in\N$, and let
$\mathcal P$ be a set of positive integers such that
$\card\{n\in\mathcal P:n\ge n_1,\ X<n\le2X\}\le C_{\mathrm w}X^{\vartheta}$
for all $X\ge X_0$. Then $\card\{n\in\mathcal P:n\le N\}\le KN^{\vartheta}$
for all $N\ge1$, with
$K=\max(2,2X_0,2n_1)+C_{\mathrm w}/(2^{\vartheta}-1)+1$.
\end{lemma}
\begin{proof}
By strong induction on the integer $m=\lfloor N\rfloor$, the assertion
is trivial for $m\le\max(2,2X_0,2n_1)$. Otherwise the prefix up to
$\lfloor m/2\rfloor$ and the window $(m/2,m]$ contribute at most
$(K+C_{\mathrm w})(m/2)^\vartheta\le Km^\vartheta$, since
$C_{\mathrm w}\le(2^\vartheta-1)K$.
\end{proof}

\section{The easy bound}\label{sec:easy}
This section bounds the number of points of a box by $\min(D,E)$, up to
$X^\delta$. From the literature it uses only $\ST$, for binary cubic forms.

\subsection{The two fibrations}
Fix a chart and put $\gamma_\pm=\sigma_\pm(\gamma_\chi)$,
$\mathsf Y_\pm(y)=(y_0\pm\sqrt\Delta y_1)/m_{\OO}$ and
$\mathsf X_\pm(x)=(x_0\pm\sqrt\Delta x_1)/m_{\OO}$.
Then \eqref{eq:coeffid} gives
\begin{equation}\label{eq:presentation}
 \frac{2\sqrt\Delta}{e_\chi}\,F_\chi(y;x)
 =\gamma_+\mathsf Y_+(y)^3\mathsf X_+(x)^2-\gamma_-\mathsf Y_-(y)^3\mathsf X_-(x)^2 .
\end{equation}
Here $\gamma_+\gamma_-\ne0$; each pair of linear forms is independent,
and its values at nonzero integer vectors are nonzero, since their
product is the corresponding norm.

\begin{lemma}\label{lem:cubes}
Let $k_1,k_2\in\C$ be nonzero and let $P_1,P_2\in\C[t_0,t_1]$ be linearly
independent linear forms. Then $k_1P_1^3-k_2P_2^3$ is nonzero and squarefree.
\end{lemma}
\begin{proof}
After rescaling $P_1,P_2$ by cube roots of $k_1,k_2$, the form becomes
$P_1^3-P_2^3$. Its three linear factors are distinct and nonzero because
$P_1,P_2$ are independent.
\end{proof}

\begin{lemma}[Cubic fibres]\label{lem:cubicfibre}
For every chart $\chi$ and every $x\ne0$ in $\Z^2$ the number
of primitive $y\in\Z^2$ with $F_\chi(y;x)=e_\chi$ is at most
$K_{\mathrm c}(\chi)=5600\,\tau(e_\chi)\,3^{1+\omega(e_\chi)}$.
\end{lemma}
\begin{proof}
By \eqref{eq:presentation} and Lemma~\ref{lem:cubes},
$F_\chi(\cdot\,;x)$ is a nonzero squarefree cubic. If solutions exist,
its content $g$ divides $e_\chi$ and $F_\chi(\cdot\,;x)=gQ_x$, with
$Q_x$ primitive and squarefree. For each of the $2\tau(e_\chi)$ signed
divisors $m$ of $e_\chi$, hypothesis $\ST$ gives at most
$2800\cdot3^{1+\omega(e_\chi)}$ primitive solutions of $Q_x=m$.
This bounds in particular the level $m=e_\chi/g$.
\end{proof}

\begin{lemma}[Quadratic fibres]\label{lem:quadfibre}
Let $h\ne0$ be an integer and $\eps>0$. There is $K_{\mathrm q}=K_{\mathrm q}(h,\eps)$ such that for
all real $H_1\ge1$, $Z_1\ge2$ and all integers $a_1,a_2,a_3$ of absolute
value at most $H_1$ with $a_2^2-4a_1a_3\ne0$, the number of primitive
$(u_1,u_2)\in\Z^2$ with $|u_1|,|u_2|\le Z_1$ and
$a_1u_1^2+a_2u_1u_2+a_3u_2^2=h$ is at most $K_{\mathrm q}(H_1Z_1)^{\eps}$.
\end{lemma}
\begin{proof}
Write $\psi$ for the form and $D_1=a_2^2-4a_1a_3$. In each of the
$4h^2$ residue classes modulo $2h$ containing solutions, fix one $w$.
For another solution $u$ in that class put
\[
 P=2a_1u_1w_1+a_2(u_1w_2+u_2w_1)+2a_3u_2w_2,
 \qquad R=u_1w_2-u_2w_1.
\]
Expansion gives $P^2-D_1R^2=4h^2$. Both $P-2h$ and $R$ are divisible
by $2h$, so $(p,r)=(P/(2h),R/(2h))$ is integral, satisfies
$p^2-D_1r^2=1$, and has $|p|,|r|\le K_1=\lceil6H_1Z_1^2\rceil$.
The map is injective because
$w_1P+(a_2w_1+2a_3w_2)R=2hu_1$ and
$w_2P-(2a_1w_1+a_2w_2)R=2hu_2$.
For $D_1<0$ there are at most four such Pell solutions, and for square
$D_1$ at most two. Otherwise they are $\pm\epsilon_1^j$ for a
fundamental Pell unit with rational coordinate at least $2$;
the rational coordinate of $\epsilon_1^j$ is at least $2^{|j|}$.
Thus there are at most $4\log K_1/\log2+4$ per class, and
$4h^2(4\log K_1/\log2+4)\ll_{h,\eps}(H_1Z_1)^\eps$.
\end{proof}

\begin{lemma}\label{lem:quaddisc}
For every chart $\chi$ and $y\ne0$ in $\Z^2$, the binary quadratic form
$x\mapsto F_\chi(y;x)$ has nonzero discriminant.
\end{lemma}
\begin{proof}
Write $\mathsf X_\pm(x)=r_\pm x_0+s_\pm x_1$ and $k_1=\gamma_+\mathsf Y_+(y)^3$,
$k_2=\gamma_-\mathsf Y_-(y)^3$, which are nonzero. By \eqref{eq:presentation} the
form is a nonzero multiple of $k_1\mathsf X_+^2-k_2\mathsf X_-^2$, whose discriminant is
$4k_1k_2(r_+s_--r_-s_+)^2\ne0$.
\end{proof}

\subsection{The bound}
\begin{proposition}[Easy bound]\label{prop:easy}
Let $H$ be as in Proposition~\ref{prop:charts}. For every
$\delta>0$ there is $K_{\mathrm e}$ such that for all $X\ge2$ and
$D,E\ge1$ with $D^3E^2\le128H^6X^2$, the number of triples
$(\chi,y,x)\in\mathcal S(X)$ with $U\le D$ and $V\le E$ is at most
$K_{\mathrm e}X^{\delta}\min(D,E)$.
\end{proposition}
\begin{proof}
For a fixed chart, $\hgt(y)^2\le D$, $\hgt(x)^2\le E$, and both
heights are at most $12H^3X$. There are at most $9E$ possible $x$;
Lemma~\ref{lem:cubicfibre} bounds each fibre by $K_{\mathrm c}(\chi)$.
There are at most $9D$ possible $y$. In each such primitive fibre,
Lemma~\ref{lem:quaddisc} gives nonzero discriminant, the quadratic
coefficients are at most $4\|F_\chi\|(12H^3X)^3$, and the variables
are bounded by $12H^3X$. Lemma~\ref{lem:quadfibre} with
$h=e_\chi$, $\eps=\delta/4$ bounds this fibre by $K_\chi X^\delta$.
Taking the smaller bound and summing over the finite chart list proves
the proposition.
\end{proof}
\section{The determinant method on a cell}\label{sec:det}
We construct an auxiliary form of bidegree $(3q,2q)$ vanishing on the
points near a short branch segment of $R_\phi=0$. This uses the chart
bounds, with no hypothesis from the literature.

\subsection{Branches and cells}
\begin{proposition}[Branches]\label{prop:branches}
For every frame $\phi$ there are an integer $J\ge0$, intervals
$[a_j,b_j]\subset[-2,2]$ and functions $g_j$ $(1\le j\le J)$, and constants
$r_0,c_{\mathrm g},C_{\mathrm g},\delta_0>0$, such that:
\begin{itemize}
\item[(i)] $g_j$ is smooth on $(a_j-2r_0,b_j+2r_0)$, and for
$\xi\in[a_j-r_0,b_j+r_0]$ we have $R_\phi(\xi,g_j(\xi))=0$ and
$|g_j'(\xi)|\ge c_{\mathrm g}$;
\item[(ii)] for every $\varkappa\ge0$ there is $M_\varkappa$ with
$|g_j^{(i)}(\xi)|\le M_\varkappa$ for all $j$, all $i\le \varkappa+1$ and all
$\xi\in[a_j-r_0,b_j+r_0]$;
\item[(iii)] every $(\xi,\zeta)$ in the balanced patch of $\phi$ with
$|R_\phi(\xi,\zeta)|\le\delta_0$ has an index $j$ with $\xi\in[a_j,b_j]$ and
$|\zeta-g_j(\xi)|\le C_{\mathrm g}|R_\phi(\xi,\zeta)|$.
\end{itemize}
\end{proposition}
\begin{proof}
Write $W_\pm=\sigma_\pm(\gamma_\chi)\ell_\pm^3m_\pm^2$.
At a zero of $R_\phi$ in the balanced patch, $W_+=W_-\ne0$, and
\[
 \ell_+\ell_-\,\partial_\xi(W_+-W_-)=3W_+(\ell_+'\ell_--\ell_-'\ell_+),\qquad
 m_+m_-\,\partial_\zeta(W_+-W_-)=2W_+(m_+'m_--m_-'m_+).
\]
The brackets are nonzero constants by independence of the linear
forms; thus both partial derivatives of $R_\phi$ are nonzero.
At such a point $(\xi_0,\zeta_0)$ write
$R_\phi(\xi,\zeta_0+t)=\varphi_2(\xi)t^2+\varphi_1(\xi)t+\varphi_0(\xi)$.
Here $\varphi_0(\xi_0)=0$ and
$\varphi_1(\xi_0)\varphi_0'(\xi_0)\ne0$. Locally the smooth function
\[
 g(\xi)=\zeta_0-
 \frac{2\varphi_0(\xi)}{\varphi_1(\xi)+\operatorname{sgn}(\varphi_1(\xi_0))
 \sqrt{\varphi_1(\xi)^2-4\varphi_2(\xi)\varphi_0(\xi)}}
\]
satisfies $R_\phi(\xi,g(\xi))=0$ and
$g'(\xi_0)=-\varphi_0'(\xi_0)/\varphi_1(\xi_0)\ne0$.
The factorization
\[
 R_\phi(\xi,\zeta)=(\zeta-g(\xi))
 [\varphi_1(\xi)+\varphi_2(\xi)(\zeta+g(\xi)-2\zeta_0)]
\]
gives $|\zeta-g(\xi)|\le2|\varphi_1(\xi_0)|^{-1}|R_\phi(\xi,\zeta)|$
in a sufficiently small box.

For each zero $z=(s_z,w_z)$ choose such an open box $B_z$ with horizontal
interval $(s_z-h_z,s_z+h_z)$, where $0<h_z\le1/2$, so that $g_z$ is
smooth on $(s_z-4h_z,s_z+4h_z)$ and its equation and
$|g_z'|\ge m_z>0$ hold on $[s_z-2h_z,s_z+2h_z]$.
Let $C_z$ be the residual constant on $B_z$.
The zero set in the compact balanced patch $S$ has a finite subcover
$B_1,\ldots,B_J$. Put $a_j=s_j-h_j$, $b_j=s_j+h_j$,
$r_0=\frac12\min_jh_j$, $c_{\mathrm g}=\min_jm_j$ and
$C_{\mathrm g}=1+\sum_jC_j$. Continuity on the compact intervals
$[s_j-2h_j,s_j+2h_j]$ bounds all derivatives of each prescribed finite
order, proving (i)--(ii). The compact set $S\setminus\bigcup_jB_j$
contains no zero; take $\delta_0$ to be half its positive minimum of
$|R_\phi|$, or $1$ if this set is empty. This proves (iii).
If the zero set is empty, take $J=0$ and arbitrary positive graph
constants, and choose $\delta_0$ by the same minimum on $S$ (or $1$
if $S$ is empty).
\end{proof}
The number of branches and the constants come from a compactness argument
and are not effective.

For an integer $M\ge1$ and a branch $j$ put $\xi_{j,k}=a_j+k/M$ for
$0\le k\le(b_j-a_j)M$. The intervals $[\xi_{j,k},\min(b_j,\xi_{j,k}+1/M)]$
cover $[a_j,b_j]$, there are at most $5M$ of them, the points $\xi_{j,k}$
are $1/M$-separated, and by the mean value theorem and (i) the values
$g_j(\xi_{j,k})$ are $c_{\mathrm g}/M$-separated.

\begin{definition}\label{def:cell}
Let $57/305<\mu<1/5$, $X\ge2$ and $M=\lceil X^\mu\rceil$. A \emph{cell}
$\mathcal C$ consists of a frame $\phi$, a branch $j$ and an index $k$ as
above; $\xi_{\mathcal C}=\xi_{j,k}$. Its set of points
$\mathcal S_{\mathcal C}(X)$ consists of the triples
$(\chi,y,x)\in\mathcal S(X)$ for which the frame given by some choice of
dominant coordinates is $\phi$ and the slopes satisfy
$\xi_{\mathcal C}\le\xi\le\min(b_j,\xi_{\mathcal C}+1/M)$ and
$|\zeta-g_j(\xi)|\le C_{\mathrm g}|R_\phi(\xi,\zeta)|$. For such a point we
write $p=(p_0,p_1,p_2,p_3)$ for the coordinates of $y$ and $x$ ordered so
that $p_1=d_y$ and $p_3=d_x$ are the dominant ones; thus $\xi=p_0/p_1$ and
$\zeta=p_2/p_3$.
\end{definition}
If $X\ge H^3/\delta_0$, Lemma~\ref{lem:slopes} and (iii) show that every
triple of $\mathcal S(X)$ lies in $\mathcal S_{\mathcal C}(X)$ for at least
one cell.

\subsection{Weighted jets}
Let $0<\mu<1/5$. In $\R[u,v]$ give the monomial $u^iv^j$ the \emph{weight}
$\mu i+j$, and say that a polynomial has weighted order at least $t$ if all
its monomials have weight at least $t$; weighted orders add under
multiplication. Let $J_\mu(t)=\card\{(i,j)\in\N^2:\mu i+j<t\}$. Fix a
polynomial $g\in\R[u]$, and for $P\in\R[u,v]$ put $P^g(u,v)=P(u,g(u)+v)$.
For integers $a,b\ge0$ and real $t$ let $W^{a,b}_t(g)$ be the space of
$P$ with $\deg_uP\le a$, $\deg_vP\le b$ such that $P^g$ has weighted order
at least $t$, and $W_{q,t}(g)=W^{3q,2q}_t(g)$. Since $W^{a,b}_t(g)$ is the
kernel of the linear map sending $P$ to the coefficients of $P^g$ at the
monomials of weight less than $t$,
\begin{equation}\label{eq:codim}
 \dim W^{a,b}_t(g)\ge(a+1)(b+1)-J_\mu(t).
\end{equation}

\begin{lemma}[Contact polynomials]\label{lem:contact}
For $(a,b)=(2,1)$ and for $(a,b)=(1,2)$ there is a nonzero $P\in\R[u,v]$
with $\deg_uP\le a$, $\deg_vP\le b$ such that $P^g$ has weighted order at
least $5\mu$.
\end{lemma}
\begin{proof}
The six-dimensional space of such $P$ contains a nonzero solution to
the five conditions that the coefficients of $u^0,\ldots,u^4$ in
$P(u,g(u))$ vanish. The pure $u$ terms of $P^g$ then have weight at
least $5\mu$, and every term involving $v$ has weight at least $1>5\mu$.
\end{proof}
We fix such polynomials $P_{21}$ and $P_{12}$.

\begin{lemma}[Mixed-contact bound]\label{lem:mixed}
Let $q,k,l\ge0$ be integers with $2k+l\le3q$ and $k+2l\le2q$. For every real
$t$,
\[
 \dim W_{q,t}(g)\ge(3q-2k-l+1)(2q-k-2l+1)-J_\mu\bigl(t-5\mu(k+l)\bigr).
\]
\end{lemma}
\begin{proof}
Multiplication by the nonzero $P_{21}^kP_{12}^l$, whose shift has
weighted order at least $5\mu(k+l)$, injects
$W^{3q-2k-l,\,2q-k-2l}_{t-5\mu(k+l)}(g)$ into $W_{q,t}(g)$.
Apply \eqref{eq:codim} to the source.
\end{proof}

\begin{lemma}\label{lem:triangle}
For $h\ge0$, $J_\mu(h)\le h^2/(2\mu)+(1+1/\mu)h+1$.
\end{lemma}
\begin{proof}
Let $n=\lfloor h\rfloor$. For each $0\le j\le n$ there are at most
$(h-j)/\mu+1$ values of $i$, so
$J_\mu(h)\le((n+1)h-n(n+1)/2)/\mu+n+1$, and the claimed bound exceeds this
by $((h-n)^2+n)/(2\mu)+(h-n)\ge0$.
\end{proof}

\begin{lemma}[Layer-cake inequality]\label{lem:layer}
Let $(C_{ij})$ be a real matrix with rows $i\in I$, $I$ finite, and columns
$1\le j\le n$, let $w_i\ge0$ be weights, and let $\operatorname{nul}(t)$ be the dimension
of the space of $x\in\R^n$ with $\sum_jC_{ij}x_j=0$ for all $i$ with
$w_i<t$. If the rows $i_1,\dots,i_n$ form a nonsingular $n\times n$ matrix,
then $\int_0^\infty\operatorname{nul}(t)\,dt\le w_{i_1}+\dots+w_{i_n}$.
\end{lemma}
\begin{proof}
On that space the map $x\mapsto(\sum_jC_{i_lj}x_j)_{l:\,w_{i_l}\ge t}$ is
injective, so $\operatorname{nul}(t)\le\card\{l:w_{i_l}\ge t\}$; integrate over $t$.
\end{proof}

\begin{proposition}[Weight surplus]\label{prop:surplus}
Let $n_q=(3q+1)(2q+1)$, the number of monomials $u^av^b$ with $a\le3q$,
$b\le2q$. Let $57/305<\mu<1/5$, $\rho(\mu)=(305\mu-57)/18$ and
$b_\mu=25\mu/3$, and let
$q\ge1$ be an integer with
\begin{equation}\label{eq:qlarge}
 q>\frac{2}{\rho(\mu)}\Bigl(b_\mu^2\bigl(11+\tfrac1\mu\bigr)+b_\mu(60\mu+11)+6\Bigr).
\end{equation}
Let $g\in\R[u]$, let $(\mathsf b_j)$ be a basis of the space of polynomials of degree
at most $3q$ in $u$ and $2q$ in $v$, and let $m_1,\dots,m_{n_q}$ be
monomials in $u,v$ such that the matrix of the coefficients of $m_i$ in
$\mathsf b_j^g$ is nonsingular. Then the sum of the weights of the $m_i$ is greater
than $n_qq+\rho(\mu)q^3/2$.
\end{proposition}
\begin{proof}
Apply Lemma~\ref{lem:layer} to this matrix, with $w_i$ the weight of $m_i$.
If $P=\sum_jx_j\mathsf b_j$ lies in $W_{q,t}(g)$ then $x$ lies in the space of the
lemma, so $\dim W_{q,t}(g)\le\operatorname{nul}(t)$ and it suffices to bound
$\int_0^{qb_\mu}\dim W_{q,t}(g)\,dt$ from below.

To select the contacts, write $r=s_1+s_2$ under the constraints
$s_1,s_2\ge0$, $2s_1+s_2\le3$, $s_1+2s_2\le2$.
Then $0\le r\le5/3$ and the residual area is
$(3-r-s_1)(2-2r+s_1)$, with derivative $1+r-2s_1$.
For fixed $r$ it is maximized at $s_1=r$ if $r\le1$, and at
$s_1=(1+r)/2$ otherwise. Its maximum is
\[
 A(r)=\begin{cases}6-7r+2r^2,&0\le r\le1,\\
 (5-3r)^2/4,&1\le r\le5/3.\end{cases}
\]
At level $z$, maximize $A(r)-(z-5\mu r)_+^2/(2\mu)$.
As $A$ decreases, one may restrict to $r\le z/(5\mu)$.
On the two ranges its derivatives are respectively
\[
 5z-7-(25\mu-4)r,\qquad
 \tfrac12(10z-15-(50\mu-9)r).
\]
They decrease and agree at $r=1$. Their zero, clipped at $0$, gives
the following path, with transitions at $z=7/5$, $5\mu+3/5$ and
endpoint $r(b_\mu)=5/3$.

For $0\le z\le b_\mu$ define
\[
 r(z)=\begin{cases}0,&z\le\frac75,\\[2pt]
 \dfrac{5z-7}{25\mu-4},&\frac75<z\le5\mu+\frac35,\\[8pt]
 \dfrac{10z-15}{50\mu-9},&5\mu+\frac35<z\le b_\mu,\end{cases}
 \qquad
 s_1(z)=\begin{cases}r(z),&r(z)\le1,\\ \frac{1+r(z)}2,&r(z)>1,\end{cases}
\]
and $s_2(z)=r(z)-s_1(z)$. Then $s_1,s_2\ge0$, $2s_1+s_2\le3$, $s_1+2s_2\le2$,
$0\le z-5\mu r\le b_\mu$, and the function
$\Pi(z)=(3-2s_1-s_2)(2-s_1-2s_2)-(z-5\mu r)^2/(2\mu)$ equals
\[
 6-\frac{z^2}{2\mu},\qquad
 6-\frac{z^2}{2\mu}+\frac{(5z-7)^2}{2(25\mu-4)},\qquad
 \frac{9(b_\mu-z)^2}{2\mu(50\mu-9)}
\]
on the three ranges (here $25\mu-4>0$ and $50\mu-9>0$).
Take $k=\lfloor qs_1\rfloor$ and $l=\lfloor qs_2\rfloor$ in
Lemma~\ref{lem:mixed} with $t=qz$. The two sides of the rectangle are at
least $q(3-2s_1-s_2)$ and $q(2-s_1-2s_2)$, and the level
$h=qz-5\mu(k+l)$ satisfies $0\le h\le q(z-5\mu r)+10\mu$. By
Lemma~\ref{lem:triangle},
\[
 J_\mu(h)\le\frac{q^2(z-5\mu r)^2}{2\mu}+q\,b_\mu\Bigl(11+\frac1\mu\Bigr)+60\mu+11,
\]
so $\dim W_{q,qz}(g)\ge q^2\Pi(z)-qb_\mu(11+1/\mu)-(60\mu+11)$. Integrating the three displayed quadratic pieces gives
$\int_0^{b_\mu}\Pi(z)\,dz=17/6+305\mu/18=6+\rho(\mu)$.
Substituting $t=qz$,
\[
 \sum_iw_i\ge\int_0^{qb_\mu}\dim W_{q,t}(g)\,dt
 \ge q^3\bigl(6+\rho(\mu)\bigr)-q^2b_\mu^2\Bigl(11+\frac1\mu\Bigr)-q\,b_\mu(60\mu+11).
\]
Since $n_qq=6q^3+5q^2+q$, the right side exceeds $n_qq+\rho(\mu)q^3/2$ by
\eqref{eq:qlarge}.
\end{proof}
Here $6$ is the leading denominator-clearing cost in $n_qq$ and
$57/305$ is the zero of $\rho$. The path optimizes the continuous bound
for these two contact polynomials; this asserts no optimality among
other determinant constructions. The bound holds for every $g$.

\subsection{Vanishing of the determinant}
To a point $p$ of a cell attach the row
$\bigl(p_0^ap_1^{3q-a}p_2^bp_3^{2q-b}\bigr)_{0\le a\le3q,\ 0\le b\le2q}\in\Z^{n_q}$.
\begin{lemma}[Row cap]\label{lem:rowcap}
For a point of a cell, every entry of its row has absolute value at most
$(2H^3X)^q$.
\end{lemma}
\begin{proof}
The square of an entry is at most
$\hgt(y)^{6q}\hgt(x)^{4q}\le(HD_y)^{3q}(HD_x)^{2q}=(H^5D_y^3D_x^2)^q
\le(H^6n^2)^q\le(2H^3X)^{2q}$, by Definition~\ref{def:point}.
\end{proof}

\begin{lemma}[Uniform Taylor estimate]\label{lem:taylor}
Let $\varkappa,q\ge0$ be integers and $r_0,M_\varkappa>0$. There is $K_{\mathrm T}$ such
that for every function $g$ of class $C^{\varkappa+1}$ on $(-r_0,r_0)$ whose
derivatives of order at most $\varkappa+1$ are bounded by $M_\varkappa$, every real $c_0$
with $|c_0|\le2$ and all $a\le3q$, $b\le2q$, the polynomial
$P_{ab}(u,v)=(c_0+u)^a(T_\varkappa g(u)+v)^b$, where $T_\varkappa g$ is the Taylor polynomial
of $g$ at $0$ of degree $\varkappa$, has all coefficients at most $K_{\mathrm T}$ in
absolute value and satisfies
$|(c_0+u)^a(g(u)+v)^b-P_{ab}(u,v)|\le K_{\mathrm T}|u|^{\varkappa+1}$ for
$|u|\le\min(r_0/2,1)$ and $|v|\le1$.
\end{lemma}
\begin{proof}
Put $M=M_\varkappa$, $T=T_\varkappa g$ and
$S=1+(\varkappa+1)M$. Each coefficient of $T$ is bounded by $M$;
submultiplicativity of the coefficient $\ell^1$ norm therefore gives
$\|P_{ab}\|_1\le3^{3q}S^{2q}$. Taylor's theorem on the segment from
$0$ to $u$ gives $|g(u)-T(u)|\le M|u|^{\varkappa+1}/(\varkappa+1)!$.
Since $|g(u)+v|,|T(u)+v|\le S$ and $|c_0+u|\le3$, the identity
$s^b-t^b=(s-t)\sum_{j=0}^{b-1}s^{b-1-j}t^j$ bounds the difference
by $3^{3q}(2q)S^{2q}M|u|^{\varkappa+1}$ for $b\ge1$; for $b=0$
it is zero. Thus $K_{\mathrm T}=1+3^{3q}S^{2q}(1+2qM)$ works uniformly.
\end{proof}

\begin{proposition}[Vanishing]\label{prop:vanish}
Let $57/305<\mu<1/5$ and let $q\ge1$ be an integer satisfying
\eqref{eq:qlarge}, chosen in terms of $\mu$ only. There is $X_1\ge2$ such that for all $X\ge X_1$, every
cell $\mathcal C$ and all points $p^{(1)},\dots,p^{(n_q)}$ of
$\mathcal S_{\mathcal C}(X)$, the $n_q\times n_q$ integer matrix formed by
their rows has determinant zero.
\end{proposition}
\begin{proof}
Put $c_0=\xi_{\mathcal C}$, $g(u)=g_j(c_0+u)$,
$u_i=\xi_i-c_0$ and $v_i=\zeta_i-g(u_i)$. Then
$0\le u_i\le X^{-\mu}$ and $|v_i|\le C_{\mathrm g}H^3/X$.
Set $\delta_1=\rho(\mu)q^3/4$ and choose $\varkappa$ so that
$\mu(\varkappa+1)>n_qq+\delta_1+1$.
Lemma~\ref{lem:taylor}, with the uniform branch bounds, approximates
the normalized matrix $\mathcal E_{i,(a,b)}=\xi_i^a\zeta_i^b$ by
\[
 \mathcal E=\mathcal V\mathcal K+\mathcal E',\qquad
 \mathcal V_{i,\epsilon}=u_i^{\epsilon_1}v_i^{\epsilon_2},\quad
 \epsilon_1+\epsilon_2\le\varkappa,\qquad
 |\mathcal E'_{i,(a,b)}|\ll X^{-\mu(\varkappa+1)}.
\]
Here $\mathcal K$ is the bounded coefficient matrix of
$P_{ab}=(c_0+u)^a(T_\varkappa g(u)+v)^b$ restricted to those monomials.
The omitted monomials have weight at least $\mu(\varkappa+1)$,
which justifies their inclusion in the error.
Expand $\det(\mathcal V\mathcal K)$ by multilinearity in its rows.
Each term is a coefficient determinant times
$\prod_i u_i^{\epsilon_{i1}}v_i^{\epsilon_{i2}}$.
If the coefficient determinant is nonzero,
Proposition~\ref{prop:surplus}, applied to $T_\varkappa g$ and the basis
$(c_0+u)^av^b$, makes the total weight greater than
$n_qq+2\delta_1$. The finitely many terms thus give
$|\det(\mathcal V\mathcal K)|\ll X^{-n_qq-2\delta_1}$.
All matrix entries are bounded; terms containing an error row give
$O(X^{-\mu(\varkappa+1)})$. Consequently
$|\det\mathcal E|\ll X^{-n_qq-\delta_1}$.
Clearing row $i$ by $p_1^{3q}p_3^{2q}$ multiplies the determinant by
at most $(2H^3X)^{qn_q}$ (Lemma~\ref{lem:rowcap}). The integer
determinant is therefore $O(X^{-\delta_1})$, uniformly over the
finite branch list, and is zero for sufficiently large $X$.
\end{proof}

\begin{proposition}[Auxiliary form]\label{prop:aux}
Let $57/305<\mu<1/5$ and $q$ be as in Proposition~\ref{prop:vanish}. There is
$X_{\mathrm a}\ge2$ such that for every $X\ge X_{\mathrm a}$ and every cell
$\mathcal C$ there is a nonzero form
$G_{\mathcal C}\in\Z[p_0,p_1;p_2,p_3]$, bihomogeneous of bidegree $(3q,2q)$,
with $\|G_{\mathcal C}\|\le X^{qn_q+1}$ and $G_{\mathcal C}(p)=0$ for every
point $p$ of $\mathcal S_{\mathcal C}(X)$.
\end{proposition}
\begin{proof}
Take $X_{\mathrm a}=\max(X_1,n_q!\,(2H^3)^{qn_q})$. The distinct
integer rows form a finite matrix $\mathcal M$ by
Lemma~\ref{lem:rowcap}; its rank $r<n_q$ by
Proposition~\ref{prop:vanish}. If $r=0$, take any coordinate vector
$c$; otherwise choose a nonzero $r$-minor and one further column.
The signed $r$-minors of this $r\times(r+1)$ block, extended by zero,
give $0\ne c\in\Z^{n_q}$ with $\mathcal Mc=0$: every row pairing is
a vanishing $(r+1)$-minor. Moreover
$|c_{(a,b)}|\le r!(2H^3X)^{qr}\le X^{qn_q+1}$.
The form $\sum_{a,b}c_{(a,b)}p_0^ap_1^{3q-a}p_2^bp_3^{2q-b}$
has the required properties.
\end{proof}
After the permutation of coordinates given by the frame it is a
form in $(y;x)$ vanishing at the triples of the cell.
\section{Points on the auxiliary curve}\label{sec:components}
We count $F_\chi=e_\chi$ on an auxiliary curve using the fibre bounds,
the three literature hypotheses, and the lifted-curve count of
\cite{powerful}. For a chart $\chi$ and $U_y,U_x\in\mathrm{GL}_2(\Z)$ put
\[
 F(\lambda;\tau)=F_\chi(U_y\lambda;U_x\tau),\qquad\lambda,\tau\in\Z^2 .
\]
This is again an integral form of bidegree $(3,2)$, and by
\eqref{eq:presentation}
\begin{equation}\label{eq:presentation2}
 \frac{2\sqrt\Delta}{e_\chi}\,F(\lambda;\tau)
 =\gamma_+\mathcal L_+(\lambda)^3\mathcal M_+(\tau)^2-\gamma_-\mathcal L_-(\lambda)^3\mathcal M_-(\tau)^2
\end{equation}
with $\gamma_\pm\ne0$ and linear forms $\mathcal L_\pm=\mathsf Y_\pm\circ U_y$,
$\mathcal M_\pm=\mathsf X_\pm\circ U_x$ over $\C$; the forms $\mathcal L_+,\mathcal L_-$ are linearly
independent, and so are $\mathcal M_+,\mathcal M_-$. The maps $\lambda\mapsto U_y\lambda$ and
$\tau\mapsto U_x\tau$ are injective and preserve primitivity.

\begin{proposition}[Component count]\label{prop:components}
Let $D_0,\varpi\ge1$ be integers and
$\delta>0$. There is $K$ such that the following holds for every real
$X\ge2$, every chart $\chi\in\mathcal X$, all
$U_y,U_x\in\mathrm{GL}_2(\Z)$ with entries of absolute value at most
$X^\varpi$, every nonzero $Q\in\Z[\lambda_1,\lambda_2;\tau_1,\tau_2]$,
bihomogeneous of a bidegree $(a_0,b_0)$ with $a_0+b_0\le D_0$ and with
$\|Q\|\le X^\varpi$, and all real $\Lambda_1\ge\Lambda_2>0$,
$\Upsilon_1\ge\Upsilon_2>0$ with $\Lambda_1,\Upsilon_1\le X^\varpi$: the set of
pairs of primitive vectors $(\lambda,\tau)$ with
\[
 |\lambda_i|\le\Lambda_i,\quad|\tau_i|\le\Upsilon_i\quad(i=1,2),\qquad
 F(\lambda;\tau)=e_\chi,\qquad Q(\lambda;\tau)=0
\]
is finite and has at most
$K X^{\delta}\bigl(1+(\Lambda_1^3\Upsilon_1^2)^{1/10}\bigr)$ elements.
\end{proposition}
Call the pairs in the displayed set \emph{solutions}; $K$ may depend on
$(A,B,C,D_0,\varpi,\delta)$.

\subsection{Components}
\begin{lemma}[Classification]\label{lem:factor}
Let $c_4$ be the constant in $\GE$ for four variables. Let $Q$ be as
in Proposition~\ref{prop:components}. There are at most $a_0+b_0$
polynomials $G\in\Z[\lambda;\tau]$, each irreducible, dividing $Q$ and
bihomogeneous of some bidegree $(\alpha,\beta)$ with $\alpha\le a_0$,
$\beta\le b_0$, such that every integral zero of $Q$ is a zero of one of
them and $\|G\|\le\exp(c_4(a_0+b_0))\,\|Q\|$. Each $G$ is irreducible over
$\Q$ and is of exactly one of the following types.
\begin{itemize}
\item[(a)] \emph{Pure}: $\beta=0$ and $\alpha\ge2$, or $\alpha=0$ and
$\beta\ge2$. Then $G$ has no integral zero with $\lambda\ne0$, respectively
$\tau\ne0$.
\item[(b)] \emph{Lines}: $G=p\lambda_1+p'\lambda_2$, or
$G=p\tau_1+p'\tau_2$.
\item[(c)] \emph{Low mixed}: $\alpha,\beta\ge1$ and $2\alpha+3\beta<10$,
that is, $(\alpha,\beta)\in\{(1,1),(1,2),(2,1),(3,1)\}$. Then $G$ is
irreducible over $\C$.
\item[(d)] \emph{High mixed}: $\alpha,\beta\ge1$ and $2\alpha+3\beta\ge10$.
\end{itemize}
\end{lemma}
\begin{proof}
Take the nonconstant primitive irreducible factors of $Q$.
Their total degrees sum to at most $a_0+b_0$; they are bihomogeneous,
and $\GE$ bounds their heights. Gauss's lemma gives irreducibility
over $\Q$. A binary form irreducible over $\Q$ of degree at least
two has no nonzero rational zero, giving (a). If
$\min(\alpha,\beta)=1$, the form is linear in one pair with coprime
binary coefficients; these remain coprime over $\C$, so the form
is absolutely irreducible. The positive bidegrees with
$2\alpha+3\beta<10$ are exactly the four in (c); all remaining
mixed factors belong to (d).
\end{proof}
Hypothesis $\GE$ enters the proof of Theorem~\ref{thm:fixed} only here, and its height
bound is used only for the components of type (c). Solutions are primitive,
so a component of type (a) carries none.

\begin{lemma}[Lines]\label{lem:lines}
In the setting of Proposition~\ref{prop:components}, a
component of type \textup{(b)} carries at most $K_1X^{\delta}$ solutions,
with $K_1$ depending on the chart list, $\varpi$ and $\delta$.
\end{lemma}
\begin{proof}
A nonzero linear form has at most two primitive zeros, differing by
sign. If it fixes $\lambda$, then $y=U_y\lambda$ is primitive with
$\hgt(y)\le2X^{2\varpi}$, and $x=U_x\tau$ is primitive with the
same height bound. The quadratic fibre has nonzero discriminant
and coefficients at most $32\|F_\chi\|X^{6\varpi}$;
Lemma~\ref{lem:quadfibre} with $\eps=\delta/(8\varpi+1)$ gives
$\ll X^\delta$ choices. If the line fixes $\tau$ instead,
Lemma~\ref{lem:cubicfibre} gives at most $K_{\mathrm c}(\chi)$ choices
for each sign. The unimodular maps are injective.
\end{proof}

\subsection{Low mixed components}
The next proposition is the analogue of
\cite[Proposition~5.3]{powerful} for forms of the shape
\eqref{eq:presentation2}.

\begin{proposition}[Graph components]\label{prop:low}
Let $h\ne0$ be an integer, $\alpha,\beta\ge1$ integers with
$\min(\alpha,\beta)=1$, and $\eps>0$. There is $K_{\mathrm l}=K_{\mathrm l}(h,\alpha,\beta,\eps)$
with the following property. Let $F$ be an integral form of bidegree
$(3,2)$ such that $\kappa_0F=\gamma_1\mathcal L_1(\lambda)^3\mathcal M_1(\tau)^2
-\gamma_2\mathcal L_2(\lambda)^3\mathcal M_2(\tau)^2$ with nonzero
$\kappa_0,\gamma_1,\gamma_2\in\C$, independent linear forms $\mathcal L_1,\mathcal L_2$ and
independent linear forms $\mathcal M_1,\mathcal M_2$ over $\C$. Let $H_G\ge1$ and let
$G\in\Z[\lambda;\tau]$ be bihomogeneous of bidegree $(\alpha,\beta)$,
irreducible over $\C$, with $\|G\|\le H_G$. Then the number of pairs of
primitive vectors $(\lambda,\tau)$ with $G(\lambda;\tau)=0$ and
$F(\lambda;\tau)=h$ is at most $K_{\mathrm l}H_G^{\eps}$.
\end{proposition}
\begin{proof}
If $\beta=1$, write $G=\tau_1P(\lambda)+\tau_2Q(\lambda)$,
put $r=\lambda$ and set $(a_G,b_G)=(r_1,r_2)$,
$(c_G,e_G)=(-Q,P)$. If $\alpha=1$, use the analogous expression
linear in $\lambda$, with $r=\tau$, $(a_G,b_G)=(-Q,P)$ and
$(c_G,e_G)=(r_1,r_2)$. In either case these are coprime integral
pairs of degrees $\beta,\alpha$, with coefficients bounded by $H_G$,
and $r$ is primitive. Coprimality follows from absolute irreducibility.
Every counted pair satisfies
\begin{equation}\label{eq:graph}
 \lambda_1b_G(r)=\lambda_2a_G(r),\qquad
 \tau_1e_G(r)=\tau_2c_G(r).
\end{equation}
For coprime integral binary forms of degree $n$ and height at most
$H_G$, the resultant argument of \cite[proof of Proposition~5.3]{powerful}
gives a nonzero integer of size at most $((2n)!)^2H_G^{4n}$ divisible
by the gcd of their values at every primitive $r$.
Apply this to the two pairs to obtain $\mathrm{Res}_1,\mathrm{Res}_2$.
Primitivity and \eqref{eq:graph} give
$(a_G,b_G)(r)=g\lambda$, $(c_G,e_G)(r)=k\tau$, where
$0\ne g\mid\mathrm{Res}_1$ and $0\ne k\mid\mathrm{Res}_2$.
Thus the binary form $\Psi=F(a_G,b_G;c_G,e_G)$, of degree
$d_R=3\beta+2\alpha$, takes the nonzero value $g^3k^2h$ at $r$,
a divisor of $\Xi=h\mathrm{Res}_1^3\mathrm{Res}_2^2$, with
\[
 |\Xi|\le|h|((2\beta)!)^6((2\alpha)!)^4H_G^{12\beta+8\alpha}.
\]
Put $U_i=\mathcal L_i(a_G,b_G)$ and $V_i=\mathcal M_i(c_G,e_G)$.
Independence makes $(U_1,U_2)$ and $(V_1,V_2)$ coprime pairs of
degrees $\beta,\alpha$. After absorbing cube roots of $\gamma_i$,
$\kappa_0\Psi$ has the form $U_1^3V_1^2-U_2^3V_2^2$.
If there is a counted pair then $\Psi\ne0$; as
$\min(\alpha,\beta)=1$, \cite[Lemma~5.2]{powerful} gives at least
three distinct projective roots.

Let $K_\Psi$ be the product of the distinct primitive integral
irreducible factors of $\Psi$. By Gauss's lemma it has content $1$,
is squarefree of degree $3\le d_K\le d_R$, and divides $\Psi$ in
$\Z[r_1,r_2]$; hence $K_\Psi(r)\mid\Xi$ at every counted point.
For each of the $2\tau(|\Xi|)$ signed divisors $m$, hypothesis
$\ST$ gives at most
$2800d_K^{1+\omega(m)}\le2800d_R\tau(|\Xi|)^{d_R}$ primitive
solutions of $K_\Psi(r)=m$.
Each $r$ determines one vector of the counted pair and the other up
to sign by \eqref{eq:graph}. Thus there are at most
$11200d_R\tau(|\Xi|)^{d_R+1}$ pairs; the divisor bound gives
$K_{\mathrm l}H_G^\eps$.
\end{proof}

The proof also gives the uniform version
\begin{equation}\label{eq:lowuniformrhs}
 \#\{(\lambda,\tau):G=0,\ F=h,\ \lambda,\tau\text{ primitive}\}
 \ll_{\alpha,\beta,\eta} (|h|H_G)^\eta
 \qquad(\eta>0).
\end{equation}
Apply the divisor bound above with exponent
$\eta/((d_R+1)(12\beta+8\alpha))$: the resulting factor
$|h|^{\eta/(12\beta+8\alpha)}H_G^\eta$ is at most $(|h|H_G)^\eta$,
with a constant independent of $F,h$.


\begin{corollary}\label{cor:low}
In the setting of Proposition~\ref{prop:components}, a
component $G$ of type \textup{(c)} carries at most $K_2X^{\delta}$ solutions.
\end{corollary}
\begin{proof}
By \eqref{eq:presentation2}, $F$ has the shape required in
Proposition~\ref{prop:low}, $G$ is irreducible over $\C$, and
$\|G\|\le\exp(c_4D_0)X^{\varpi}$ by Lemma~\ref{lem:factor}. Apply
Proposition~\ref{prop:low} with $h=e_\chi$, $H_G=\exp(c_4D_0)X^\varpi$ and
$\eps=\delta/\varpi$, for each of the four bidegrees and each chart.
\end{proof}

\subsection{High mixed components and the lifted curve}
\begin{lemma}[Lemma P, {\cite[Lemma~5.5]{powerful}}]\label{lem:P}
Let $c_{12}$ be the constant in $\CC$ for twelve variables. Let
$F\in\Z[\lambda_1,\lambda_2;\tau_1,\tau_2]$ be bihomogeneous of bidegree
$(3,2)$, let $e$ be a nonzero integer, and let
$G\in\Q[\lambda_1,\lambda_2;\tau_1,\tau_2]$ be bihomogeneous of bidegree
$(\alpha,\beta)\ne(0,0)$ and irreducible over $\Q$. Let
$\Lambda_1\ge\Lambda_2>0$ and $\Upsilon_1\ge\Upsilon_2>0$. Then the number of
pairs of primitive vectors $(\lambda,\tau)\in\Z^2\times\Z^2$ with
$|\lambda_i|\le\Lambda_i$, $|\tau_i|\le\Upsilon_i$, $F(\lambda;\tau)=e$ and
$G(\lambda;\tau)=0$ is at most
\[
 2\,c_{12}\,d^3\Omega^{1/d}(\log\Omega+d),\qquad d=2\alpha+3\beta,\quad
 \Omega=\max(1,\Lambda_1^3\Upsilon_1^2).
\]
\end{lemma}
This is proved in \cite{powerful}: the twelve monomials of bidegree $(3,2)$
map the counted points, at most two-to-one, to integral points of height at
most $\Omega$ on an affine curve of degree $2\alpha+3\beta$ in twelve
variables, to which $\CC$ applies. The form $F$ is arbitrary, and $G$ need
not be irreducible over $\C$. We may and do assume $c_{12}\ge1$.

\begin{corollary}\label{cor:high}
In the setting of Proposition~\ref{prop:components}, a
component $G$ of type \textup{(d)} carries at most
$K_3X^{\delta}(1+(\Lambda_1^3\Upsilon_1^2)^{1/10})$ solutions.
\end{corollary}
\begin{proof}
Apply Lemma~\ref{lem:P} with $e=e_\chi$. Here $10\le d\le3D_0$,
so $\Omega^{1/d}\le1+(\Lambda_1^3\Upsilon_1^2)^{1/10}$, while
$\Omega\le X^{5\varpi}$ gives
$\log\Omega\le(5\varpi/\delta)X^\delta$.
Thus $K_3=54c_{12}D_0^3(5\varpi/\delta+3D_0)$ suffices.
\end{proof}
This is the sole use of the lifted-curve count in Theorem~\ref{thm:fixed}.
The exponent $1/10$ comes from the least remaining lifted degree,
attained by bidegree $(2,2)$.

\begin{proof}[Proof of Proposition~\ref{prop:components}]
Every solution is a zero of one of the at most $D_0$ components of
Lemma~\ref{lem:factor}. Components of type (a) carry no solution, and the
other types are bounded by Lemma~\ref{lem:lines} and
Corollaries~\ref{cor:low} and~\ref{cor:high}. Finiteness holds because the
solutions lie in a box. So the proposition holds with
$K=D_0(K_1+K_2+K_3)$.
\end{proof}
\section{Reduced coordinates, moments and the proof of Theorem~\ref{thm:fixed}}\label{sec:sum}
Throughout, $57/305<\mu<1/5$, $X\ge4$ and
$M=\lceil X^\mu\rceil$, so that $1\le M\le X$.

\subsection{Strips and reduced bases}
For $K_0\ge1$, $D>0$ and a real slope $\varsigma$, the \emph{strip}
$\Sigma(K_0,D,\varsigma)$ is the set of $v=(v_0,v_1)\in\Z^2$ with
\[
 |v_0-\varsigma v_1|\le K_0\sqrt D/M,\qquad|v_1|\le K_0\sqrt D,
\]
and the \emph{strip norm} $|v|_{\varsigma}$ of $v\in\Z^2$ is the Euclidean norm
of the vector
$\bigl(M(v_0-\varsigma v_1)/(K_0\sqrt D),\ v_1/(K_0\sqrt D)\bigr)\in\R^2$. The
strip is the set of $v$ whose image lies in the square $[-1,1]^2$.

\begin{lemma}[Strips]\label{lem:strips}
There is $K_{\mathrm s}\ge1$ such that for every cell
$\mathcal C=(\phi,j,k)$ and every point $p$ of $\mathcal S_{\mathcal C}(X)$
with $U\le D$ and $V\le E$,
\[
 (p_0,p_1)\in\Sigma(K_{\mathrm s},D,\xi_{\mathcal C}),\qquad
 (p_2,p_3)\in\Sigma(K_{\mathrm s},E,g_j(\xi_{\mathcal C})).
\]
\end{lemma}
\begin{proof}
We have $|p_1|\le\sqrt D$, $|p_3|\le\sqrt E$ and
$|p_0/p_1-\xi_{\mathcal C}|\le1/M$. If $M_{\mathrm d}$ bounds first
branch derivatives, the other slope satisfies
$|p_2/p_3-g_j(\xi_{\mathcal C})|\le C_{\mathrm g}H^3/X+M_{\mathrm d}/M
\le(C_{\mathrm g}H^3+M_{\mathrm d})/M$.
Multiply by the dominant coordinates and take
$K_{\mathrm s}=1+C_{\mathrm g}H^3+M_{\mathrm d}$, maximized over frames.
\end{proof}

\begin{proposition}[Reduced bases]\label{prop:bridge}
Let $K_0\ge1$, let $D>M$ with $M$ as above, and let $|\varsigma|\le K_0$. There are
$\mathsf B\in\mathrm{GL}_2(\Z)$ and real numbers $\Lambda_1\ge\Lambda_2>0$ with
$\Lambda_1\le2K_0\sqrt D$ such that the entries of $\mathsf B$ and of $\mathsf B^{-1}$ are
at most $4(K_0+1)D$ in absolute value, the first column $b_1$ of $\mathsf B$ is a
nonzero vector of least strip norm with $\Lambda_1=2/|b_1|_\varsigma$, and every
primitive $v\in\Sigma(K_0,D,\varsigma)$ equals $\mathsf B\lambda$ for a primitive
$\lambda\in\Z^2$ with $|\lambda_1|\le\Lambda_1$ and
$|\lambda_2|\le\Lambda_2$.
\end{proposition}
\begin{proof}
Let $b_1$ be a nonzero integer vector of least strip norm; it exists because
the vectors of bounded strip norm lie in a finite box, and it is primitive.
Choose $w\in\Z^2$ with $\det(b_1,w)=1$ and put $b_2=w-kb_1$ with $k$ the
integer nearest to $\langle b_1,w\rangle_\varsigma/|b_1|_\varsigma^2$, where
$\langle\,,\rangle_\varsigma$ is the inner product of the strip norm. With
$N_i=|b_i|_\varsigma^2$ and $P_0=\langle b_1,b_2\rangle_\varsigma$ we get $N_1\le N_2$ and
$|P_0|\le N_1/2$, and therefore, for all real $l_1,l_2$,
\begin{equation}\label{eq:gram}
 |l_1b_1+l_2b_2|_\varsigma^2\ge\tfrac12\bigl(l_1^2N_1+l_2^2N_2\bigr).
\end{equation}
Let $\mathsf B$ have columns $b_1,b_2$, and $\Lambda_i=2/\sqrt{N_i}$. If
$v=\mathsf B\lambda$ lies in the strip, its image has squared length at most $2$, so
\eqref{eq:gram} gives $\lambda_i^2N_i\le4$; and $\lambda$ is primitive when
$v$ is. Write $b_1=(a_1,c_1)$, $b_2=(a_2,c_2)$ and $R_D=K_0\sqrt D$. Then
$R_D^2N_1=M^2(a_1-\varsigma c_1)^2+c_1^2\ge1$, which gives
$\Lambda_1\le2R_D$. The images of $b_1,b_2$ span a parallelogram of area
$M/R_D^2$, so $N_1N_2-P_0^2=M^2/R_D^4$ and $3N_1N_2\le4M^2/R_D^4$; hence
$R_D^2N_2\le2M^2$. So $|c_i|\le2M$ and $|a_i-\varsigma c_i|\le2$, which gives
$|a_i|\le2+2K_0M\le4(K_0+1)D$. The entries of $\mathsf B^{-1}$ are those of $\mathsf B$ up
to sign and order.
\end{proof}
The hypothesis $D>M$ is used for the bound on the entries. For each cell
$\mathcal C=(\phi,j,k)$ we fix $\mathsf B_y(\mathcal C)$ with
$\Lambda_i(\mathcal C)$ for the strip $\Sigma(K_0,D,\xi_{\mathcal C})$, and
$\mathsf B_x(\mathcal C)$ with $\Upsilon_i(\mathcal C)$ for
$\Sigma(K_0,E,g_j(\xi_{\mathcal C}))$, where $K_0\ge K_{\mathrm s}$ also
bounds $5$, the slopes $\xi_{\mathcal C}$ and the values
$g_j(\xi_{\mathcal C})$.

\begin{lemma}[Per-cell count]\label{lem:percell}
Let $q$ be as in Proposition~\ref{prop:aux}
and $\delta>0$. There are $P_{\mathrm c}$ and $X_{\mathrm c}$ such that for
all $X\ge X_{\mathrm c}$, all $D,E>M$ with $D^3E^2\le128H^6X^2$ and every
cell $\mathcal C$, the number of points of $\mathcal S_{\mathcal C}(X)$ with
$U\le D$ and $V\le E$ is at most
$P_{\mathrm c}X^{\delta}\bigl(1+(\Lambda_1(\mathcal C)^3\Upsilon_1(\mathcal C)^2)^{1/10}\bigr)$.
\end{lemma}
\begin{proof}
Let $S_y,S_x$ be the coordinate permutations of the frame.
Lemmas~\ref{lem:strips} and Proposition~\ref{prop:bridge} give an
injective parametrization
$S_yy=\mathsf B_y\lambda$, $S_xx=\mathsf B_x\tau$ by primitive
vectors in the reduced box. Set $U_y=S_y\mathsf B_y$,
$U_x=S_x\mathsf B_x$ and
$Q(\lambda;\tau)=G_{\mathcal C}(\mathsf B_y\lambda;\mathsf B_x\tau)$.
Then $F_\chi(U_y\lambda;U_x\tau)=e_\chi$ and $Q=0$; the latter
form is nonzero of bidegree $(3q,2q)$.
As $D,E\le128H^6X^2$, the matrix entries are at most $X^3$ for
large $X$. Expanding the substitution gives
\[
 \|Q\|\le(5q+1)^42^{5q}X^{qn_q+1+15q}\le X^\varpi,
 \qquad \varpi=qn_q+15q+2.
\]
Also $\Lambda_1,\Upsilon_1\le X^2$. Proposition~\ref{prop:components}
with $D_0=5q$ and this $\varpi$ now applies.
\end{proof}

\subsection{Moments}
\begin{lemma}[Tail count]\label{lem:tail}
Let $K_0\ge1$ and $c>0$. There is $C_{\mathrm t}$ such that the following
holds for all $D>0$. Let $(\varsigma_i)_{i\in I}$ be real numbers with
$|\varsigma_i|\le K_0$ and $|\varsigma_i-\varsigma_j|\ge c/M$ for $i\ne j$, with
$\card I\le K_0M$; for each $i$ let $b^{(i)}$ be a nonzero integer vector
and $\ell_i=2/|b^{(i)}|_{\varsigma_i}$. Then for every $\Lambda>0$,
\[
 \card\{i\in I:\ell_i\ge\Lambda\}\le C_{\mathrm t}D/\Lambda^2 .
\]
\end{lemma}
\begin{proof}
For $\Lambda<8K_0\sqrt D/M$, the trivial bound $K_0M$ suffices,
since $M\ge1$ and $M^2<64K_0^2D/\Lambda^2$.
Otherwise put $Z=2K_0\sqrt D/\Lambda\le M/4$.
If $\ell_i\ge\Lambda$, its nonzero vector $(p,p')$ satisfies
$|p-\varsigma_ip'|\le Z/M$ and $|p'|\le Z$.
Here $p'\ne0$, since otherwise $|p|\le1/4$; changing signs gives
$1\le p'\le Z$. Thus there are no such indices if $Z<1$.
For each integer $q=p'\le Z$, the bound
$|p|\le K_0q+1/4$ allows at most $(2K_0+2)q$ numerators.
For each $(p,q)$, admissible slopes lie in an interval of length
$2Z/(Mq)$, which contains at most $1+2Z/(cq)$ of the separated
$\varsigma_i$. Therefore the number of indices is at most
\[
 (2K_0+2)\sum_{1\le q\le Z}q\left(1+\frac{2Z}{cq}\right)
 \ll_{K_0,c}Z^2=4K_0^2D/\Lambda^2.
\]
\end{proof}
This is a double count of pairs consisting of an integer vector in a sector
and a slope close to its direction; it uses neither primitivity nor
reducedness of the $b^{(i)}$.

\begin{lemma}[Moments]\label{lem:moments}
In the setting of Lemma~\ref{lem:tail}, for $0<\vartheta<2$ there is $K_{\mathrm m}$
with $\sum_{i\in I}\ell_i^{\vartheta}\le K_{\mathrm m}M(D/M)^{\vartheta/2}$. If moreover
$(\varsigma_i')_{i\in I}$, $E>0$ and vectors with lengths $\ell_i'$ satisfy the same
hypotheses with $E$ in place of $D$, then
\[
 \sum_{i\in I}\ell_i^{3/10}\ell_i'^{\,1/5}\le
 K_{\mathrm m}'M\Bigl(\frac DM\Bigr)^{3/20}\Bigl(\frac EM\Bigr)^{1/10}
 =K_{\mathrm m}'(D^3E^2)^{1/20}M^{3/4}.
\]
\end{lemma}
\begin{proof}
Put $\Lambda_*=4K_0\sqrt{D/M}$. Lengths at most $\Lambda_*$
contribute at most $K_0M\Lambda_*^\vartheta$.
Lemma~\ref{lem:tail} bounds the contribution from
$[2^k\Lambda_*,2^{k+1}\Lambda_*)$ by
$C_{\mathrm t}M\Lambda_*^\vartheta2^\vartheta(2^\vartheta/4)^k$.
Summing this convergent series proves the first assertion.
Cauchy--Schwarz with moments $3/5$ and $2/5$ proves the second.
\end{proof}

\begin{proposition}[Hard bound]\label{prop:hard}
Let $57/305<\mu<1/5$ and $\delta>0$. There
are $K_{\mathrm h}$ and $X_{\mathrm h}$ such that for all
$X\ge X_{\mathrm h}$ and all $D,E>M$ with $D^3E^2\le128H^6X^2$, the number
of triples $(\chi,y,x)\in\mathcal S(X)$ with $U\le D$ and $V\le E$ is at
most $K_{\mathrm h}X^{\delta}\bigl(M+X^{1/10}M^{3/4}\bigr)$.
\end{proposition}
\begin{proof}
Take $X_{\mathrm h}\ge\max(4,n_0,X_{\mathrm a},X_{\mathrm c},H^3/\delta_0)$
for every frame, so all triples are covered by cells.
On each fixed branch there are at most $5M$ cells; their two slope
families are separated by $1/M$ and $c_{\mathrm g}/M$.
Apply Lemma~\ref{lem:moments} to the first basis columns, whose
lengths are $\Lambda_1,\Upsilon_1$, and sum Lemma~\ref{lem:percell}:
the branch contributes at most
\[
 P_{\mathrm c}X^\delta
 \bigl(5M+K_{\mathrm m}'(D^3E^2)^{1/20}M^{3/4}\bigr)
 \ll X^\delta(M+X^{1/10}M^{3/4}).
\]
There are finitely many frames and branches.
\end{proof}

\subsection{The schedule}
\begin{lemma}[Schedule]\label{lem:schedule}
Let $\eps>0$ and put $\eta=\min(1/1000,\eps/4)$, $\mu=57/305+\eta$ and
$\delta=\eps/4$. For all $A_1,P>0$ there is $K_{\mathrm{sch}}$ such that for every
$X\ge2$ and every family of at most $A_1(1+\log X)^2$ boxes $(D_i,E_i)$ with
counts $z_i\ge0$ satisfying
\begin{itemize}
\item[(i)] $z_i\le PX^{\delta}\min(D_i,E_i)$ for all $i$, and
\item[(ii)] $z_i\le PX^{\delta}(M+X^{1/10}M^{3/4})$ whenever $M<D_i$ and
$M<E_i$,
\end{itemize}
we have $\sum_iz_i\le K_{\mathrm{sch}}X^{293/1220+\eps}$.
\end{lemma}
\begin{proof}
Put $\vartheta=1/10+3\mu/4=293/1220+3\eta/4\ge\mu$.
Since $M\le2X^\mu$, both cases (i)--(ii) give
$z_i\le4PX^{\delta+\vartheta}$.
The number of boxes is $\ll_\eps A_1X^{\eps/4}$, and
$3\eta/4+\delta+\eps/4\le\eps$.
\end{proof}
The boxes with
$\min(D,E)\le M$ are those in which the reduced bases of
Proposition~\ref{prop:bridge} are not available; there the easy bound
suffices.

\subsection{Proof of Theorem~\ref{thm:fixed}}
Choose $\mu,\delta$ as in Lemma~\ref{lem:schedule}, and take the chart
constants and the easy and hard bounds above. For $X\ge X_{\mathrm h}$,
Lemma~\ref{lem:recover} gives $\card\mathcal W(X)\le\card\mathcal S(X)$.
Lemmas~\ref{lem:norms} and~\ref{lem:partition} split $\mathcal S(X)$
into $O((1+\log X)^2)$ boxes with $D^3E^2\le128H^6X^2$.
Propositions~\ref{prop:easy} and~\ref{prop:hard} supply the two bounds
of Lemma~\ref{lem:schedule}, with $P=K_{\mathrm e}+K_{\mathrm h}+1$.
Hence
\[
 \card\mathcal W(X)\le K_{\mathrm w}X^{293/1220+\eps}
 \qquad(X\ge X_{\mathrm h}).
\]
By Lemma~\ref{lem:growth}, $f(n)>0$ for $n\ge n_0$; thus these are
exactly the large $n$ in the window counted by $S(A,B,C;\cdot)$.
Lemma~\ref{lem:prefix} with $\vartheta=293/1220+\eps$ now yields
$S(A,B,C;N)\le KN^{293/1220+\eps}$ for every $N\ge1$. \qed

\medskip
The constant depends on $(A,B,C,\eps)$ and is not effective: class
representatives, generators, branch neighborhoods and the thresholds
$X_{\mathrm a},X_{\mathrm c}$ were chosen existentially.
The assumption $A>0$ enters growth, and $A\ne0$ also enters
Lemmas~\ref{lem:ideal} and~\ref{lem:recover}. Nonsquareness gives the
quadratic field, finite class group, nonvanishing norms, and the Pell
unit when $\Delta>0$.

\section{The uniform bound: proof of Theorem~\ref{thm:uniform}}\label{sec:uniform}
Let $f(n)=an^2+bn+c$, where $a,b,c\in\Z$, $a\ge1$,
$\Delta=b^2-4ac$ is nonsquare, and $|a|,|b|,|c|\le L$, $L\ge2$.
Constants are absolute unless a dependence is stated. The fixed-polynomial
proof leaves its ideal representatives, unit and compactness bounds
uncontrolled in $L$, and its determinant order grows as $\mu\to57/305$.
Here we give arithmetic and graph constructions with polynomial losses
and fix $\mu_0=19/100$, $q_0=4096$. We reuse
Proposition~\ref{prop:surplus} and Lemmas~\ref{lem:P}, \ref{lem:partition},
\ref{lem:moments} and~\ref{lem:prefix}.

\subsection{The trace equation}
Write $\Delta=\Delta_0r_0^2$, with $r_0\ge1$ and $\Delta_0$ squarefree,
nonsquare and $|\Delta_0|\le5L^2$. Put
$\omega_0=(1+\sqrt{\Delta_0})/2$ if $\Delta_0\equiv1\pmod4$ and
$\omega_0=\sqrt{\Delta_0}$ otherwise. Then $\OO=\Z+\Z\omega_0$.
For $y\in\Z^2$ set $\iota_0(y)=y_0+y_1\omega_0$, and let $\sigma_\pm$
send $\sqrt{\Delta_0}$ to $\pm\sqrt{\Delta_0}\in\C$.
The integral trace form
\[
 T_0(y;x)=16\operatorname{Tr}(\iota_0(y)^3\iota_0(x)^2)
\]
has bidegree $(3,2)$. Also put $z_n=(2an+b)+r_0\sqrt{\Delta_0}$,
so $z_n\bar z_n=4af(n)$.

\begin{lemma}[Balancing by a bounded-norm integer]\label{lem:A3boundedbalance}
Let $d\ge2$ be a squarefree integer. For every $R>0$ there is
$\gamma=p+q\sqrt d\in\Z[\sqrt d]$ such that
\[
 1\le |\Nm\gamma|\le3\sqrt d,
 \qquad
 \frac{R}{6d}\le
 \left|\frac{p+q\sqrt d}{p-q\sqrt d}\right|\le6dR.
\]
The integer $\gamma$ is not required to be a unit.
\end{lemma}
\begin{proof}
Let $p_j/q_j$ be the simple continued-fraction convergents of $u=\sqrt d$,
starting with $p_0=\lfloor u\rfloor$, $q_0=1$ and $p_{-1}=1$, $q_{-1}=0$.
The recurrences give $1\le q_j\le q_{j+1}$ and
$q_{j+2}\ge q_{j+1}+q_j$, hence $q_j\to\infty$. If $u_{j+1}>a_{j+1}$
is the next complete quotient, then
\[
 u=\frac{p_ju_{j+1}+p_{j-1}}{q_ju_{j+1}+q_{j-1}},\qquad
 |p_jq_{j-1}-p_{j-1}q_j|=1,
\]
so $q_{j+1}=a_{j+1}q_j+q_{j-1}$ implies
\begin{equation}
 t_j:=|p_j-uq_j|\le q_{j+1}^{-1}.                 \label{eq:cf-error}
\end{equation}
Put $A_j=p_j+uq_j$, $B_j=|p_j^2-dq_j^2|$ and $\Lambda_j=A_j^2/B_j$.
Since $p_j\ge uq_j-1>0$ and $u$ is irrational,
\begin{equation}
 1\le B_j=A_jt_j\le(2uq_j+t_j)t_j\le2u+1\le3u. \label{eq:cf-norm}
\end{equation}
Thus $q_{j+1}\le A_j$ and $A_{j+1}\le2uq_{j+1}+1\le3uA_j$, giving
\begin{equation}
 \Lambda_{j+1}\le9u^2A_j^2
 =9u^2B_j\Lambda_j\le27u^3\Lambda_j
 \le(6d)^2\Lambda_j.                            \label{eq:cf-ratio-jump}
\end{equation}
Moreover $\Lambda_j\to\infty$ because $A_j\ge uq_j$, and
$\Lambda_0\le4d$. For $R\ge1$ choose the first $j$ with
$\Lambda_j\ge R/(6d)$. If $j=0$, then $\Lambda_j\le4d\le6dR$;
otherwise \eqref{eq:cf-ratio-jump} gives $\Lambda_j<6dR$.
Now $\gamma=p_j+q_j\sqrt d$ has embedding ratio $\Lambda_j$ and
absolute norm $B_j$. For $R<1$, apply this construction to $R^{-1}$
and conjugate $\gamma$.
\end{proof}

\begin{proposition}[Arithmetic reduction]\label{prop:A3lift}
There is an absolute constant $C_0\ge1$ such that for every $n\ge4L$ with
$f(n)$ powerful there are an integer $e$ and primitive vectors
$y,x\in\Z^2$ with
\begin{itemize}
\item[(i)] $1\le e\le C_0L^{28}$ and $T_0(y;x)=e$;
\item[(ii)] $\iota_0(y)^3\iota_0(x)^2
=\dfrac{e}{32\,\Delta_0^2r_0}\,(\sqrt{\Delta_0})^3z_n$;
\item[(iii)] $|\sigma_+(\iota_0y)|\le C_0L^4|\sigma_-(\iota_0y)|$
and $|\sigma_+(\iota_0x)|\le C_0L^6|\sigma_-(\iota_0x)|$, and the same
with $\sigma_+$ and $\sigma_-$ exchanged;
\item[(iv)] $n/(C_0L^9)\le\hgt(y)^3\hgt(x)^2\le C_0L^{40}n$.
\end{itemize}
The triple $(e,y,x)$ determines $n$.
\end{proposition}
\begin{proof}
Write $d=\Delta_0$, $s=r_0$, $\theta=\sqrt d$ and $K=\Q(\theta)$.
Parentheses denote principal fractional ideals; $\Nm$ denotes the field
norm or absolute ideal norm. Set
\[
 z=(2an+b)+s\theta,\qquad w=\theta^3z=d^2s+d(2an+b)\theta,
 \qquad h_0=d^2s.
\]
Since $\Delta=ds^2$ and $s\ge1$,
\begin{equation}
 1\le|d|\le|\Delta|\le5L^2,\qquad
 0<h_0\le\Delta^2\le25L^4.                      \label{eq:arith-1}
\end{equation}
For $n\ge4L$, the inequalities $L(n+1)\le n^2/2$ and
$n^2+n+1\le2n^2$ give
\begin{equation}
 n^2/4\le f(n)\le2Ln^2,                         \label{eq:arith-2}
\end{equation}
whence $z\ne0$ and
\begin{equation}
 \Nm w=-4ad^3f(n),\qquad
 n^2\le|\Nm w|\le1000L^8n^2.                    \label{eq:arith-3}
\end{equation}

Write $f(n)=u^3v^2$ with $u,v\ge1$, by expressing every prime exponent
at least two as $3i+2j$. A prime ideal $\mathfrak p$ occurring exactly
once in $(z)$ divides $(4a\Delta)$: if $\bar z\in\mathfrak p$, then
$2s\theta=z-\bar z\in\mathfrak p$ and $4\Delta\in\mathfrak p$;
if neither $\bar z$ nor $4a$ belongs to $\mathfrak p$, valuations in
$(z)(\bar z)=(4a)(u)^3(v)^2$ give the impossible equality $1=3i+2j$.
Let $\mathfrak R$ be the product of these prime ideals. Splitting all
remaining exponents into $3i+2j$ yields integral ideals
\[
 (z)=\mathfrak R\mathfrak I^3\mathfrak J^2,
 \qquad \mathfrak R\text{ squarefree}.
\]
For $B_0=2a|\Delta|$, the identity $B_0^2=(4a\Delta)(a\Delta)$ and
squarefreeness imply $\mathfrak R\mid(B_0)$, hence
\begin{equation}
 r:=\Nm\mathfrak R\le B_0^2\le100L^6.           \label{eq:arith-4}
\end{equation}
The order $\Z[\theta]$ has discriminant $4d$, so $|D_K|\le20L^2$.
Minkowski supplies an integral ideal $\mathfrak j$ in the class of
$\mathfrak I\mathfrak J$ with
\begin{equation}
 N:=\Nm\mathfrak j\le(4/\pi)^{r_2}\frac{2!}{2^2}\sqrt{|D_K|}
 \le2\sqrt{|D_K|}\le10L.                        \label{eq:arith-5}
\end{equation}
Choose $(\xi)\mathfrak j=\mathfrak I\mathfrak J$ and set
$\alpha=z/\xi^2$, $\beta=N^3r\xi^3/z$. These are integral:
\[
 (\alpha)=\mathfrak R\mathfrak I\mathfrak j^2,\qquad
 (\beta)=\frac{(N)^3(r)\mathfrak J}{\mathfrak R\mathfrak j^3}
 \subseteq\mathfrak J.
\]
Here $(N)\subseteq\mathfrak j$, $(r)\subseteq\mathfrak R$, because
an integral ideal contains its norm (the order of its additive quotient).
The choices give the exact identity
\begin{equation}
 \alpha^3\beta^2=N^6r^2z.                        \label{eq:arith-6}
\end{equation}

For $d<0$, set $m=1$, $\alpha_1=\alpha$, $\beta_1=\beta$; conjugate
embeddings have equal moduli. For $d>0$ we have $d\ge2$ and
$2an+b\ge7L\ge3s\sqrt d$, since $s\sqrt d\le\sqrt5L$. Thus
$Z:=|\sigma_+(z)/\sigma_-(z)|\in[1/2,2]$.
Apply Lemma~\ref{lem:A3boundedbalance} to
$R=A^{-1/2}$, $A=|\sigma_+(\alpha)/\sigma_-(\alpha)|$, obtaining
\[
 \begin{gathered}
 1\le m:=|\Nm\gamma|\le3\sqrt d\le7L,\\
 F^{-1}A^{-1/2}\le G:=|\sigma_+(\gamma)/\sigma_-(\gamma)|
 \le FA^{-1/2},\quad F=6d\le30L^2.
 \end{gathered}
\]
Put $\alpha_1=\gamma^2\alpha$, $\beta_1=\bar\gamma^3\beta$.
Then $\alpha_1^3\beta_1^2=N^6r^2m^6z$ and its embedding ratios satisfy
$F^{-2}\le A_1=G^2A\le F^2$, $A_1^3B_1^2=Z$.
Since $\sqrt2F^3\le54000L^6$, in both signs of $d$ we obtain
\begin{equation}
 (900L^4)^{-1}\le A_1\le900L^4,\qquad
 (54000L^6)^{-1}\le B_1\le54000L^6.              \label{eq:arith-7}
\end{equation}
Represent $\theta\alpha_1=\iota_0(Y)$ and $\beta_1=\iota_0(X)$ in
the integral basis. Multiplication by $\theta$ preserves embedding
ratios in absolute value. Thus, with $\kappa=N^6r^2m^6$,
\begin{equation}
 1\le\kappa\le ML^{24},\qquad
 \iota_0(Y)^3\iota_0(X)^2=\kappa w,              \label{eq:arith-8}
\end{equation}
where an absolute $M\ge54000$ bounds the ratios by $ML^4,ML^6$;
for example $M=20^6\cdot100^2\cdot7^6+100000$ suffices.

Divide $Y,X$ by their positive coordinate contents $g_Y,g_X$.
The resulting $y,x$ are primitive, preserve the embedding ratios and satisfy
\begin{equation}
 \iota_0(y)^3\iota_0(x)^2=\lambda w,\qquad
 \lambda=\kappa/q,\quad q=g_Y^3g_X^2\ge1.        \label{eq:arith-9}
\end{equation}
The trace form is integral: in the half-integral basis the five factors
have denominator dividing $2^5$, and $16\operatorname{Tr}$ multiplies
the rational coordinate by $32$. Taking traces gives the positive integer
$e=T_0(y;x)$ and $qe=32\kappa h_0$, hence
\begin{equation}
 1\le e\le800ML^{28},\qquad
 \lambda=\frac e{32h_0},\qquad
 \frac1{800L^4}\le\lambda\le ML^{24}.            \label{eq:arith-10}
\end{equation}
This proves (i)--(iii) after enlarging $C_0$.

For $H(v)=\hgt(v)$ and $A(v)=|\Nm\iota_0(v)|$, solving the two
embedding equations and using $|\theta|\ge1$, $|\omega_0|,|\bar\omega_0|\le3L$
gives
\begin{equation}
 H(v)\le2\max_\pm|\sigma_\pm(\iota_0v)|,
 \qquad A(v)\le16L^2H(v)^2.                     \label{eq:arith-11}
\end{equation}
If the embedding ratio is bounded in both directions by $Q$, then also
\begin{equation}
 H(v)^2\le4QA(v).                              \label{eq:arith-12}
\end{equation}
Put $V=H(y)^3H(x)^2$. Since $A(y)^3A(x)^2=\lambda^2|\Nm w|$,
\eqref{eq:arith-3}, \eqref{eq:arith-10}--\eqref{eq:arith-12} give
\begin{equation}
 \frac{n^2}{800^2\,16^5L^{18}}\le V^2
 \le1000\,4^5M^7L^{80}n^2.                      \label{eq:arith-13}
\end{equation}
Indeed the lower bound divides by $16^5L^{10}$ and uses
$\lambda^2\ge800^{-2}L^{-8}$; the upper bound multiplies
$4^5M^5L^{24}$, $M^2L^{48}$ and $1000L^8n^2$.
Square roots prove (iv) with an absolute $C_0$ also dominating $800M$.
Finally (ii) recovers $w$ from $(e,y,x)$, and its $\theta$-coordinate
$d(2an+b)$ determines $n$ because $ad\ne0$.
\end{proof}

\subsection{The count on a shell}
Let $\mathcal P(L,Z,\Delta_0,e)$ consist of primitive pairs $(y,x)$
with $T_0(y;x)=e$, $Z\le\hgt(y)^3\hgt(x)^2\le2Z$ and the balancing
inequalities (iii) of Proposition~\ref{prop:A3lift}. Choosing dominant
coordinates gives four \emph{orientations}, each with slopes
$(\xi,\zeta)\in[-1,1]^2$. Throughout the next three propositions,
$\Delta_0$ is squarefree and nonsquare with $|\Delta_0|\le5L^2$.

\begin{proposition}[Fallback count]\label{prop:A3fallback}
For $\delta>0$ there are $K_{\mathrm f},B_{\mathrm f}>0$ such
that for all $L,Z\ge2$, all squarefree nonsquare $\Delta_0$ with
$|\Delta_0|\le5L^2$, all $1\le e\le C_0L^{28}$ and all $D,E\ge1$ with
$\sqrt D,\sqrt E\le4Z$, the number of pairs of primitive vectors with
$|y_i|\le\sqrt D$, $|x_i|\le\sqrt E$ and $T_0(y;x)=e$ is at most
$K_{\mathrm f}L^{B_{\mathrm f}}Z^{\delta}\min(D,E)$.
\end{proposition}
\begin{proof}
Write $d=\Delta_0$, $h=e$, $P=\max(C_0,4)LZ$.
For a quadratic form $Q$ of nonzero discriminant and coefficient
height at most $A\ge1$, Lemma~\ref{lem:quadfibre} gives, uniformly in
$h\ne0$ and $X\ge2$,
\begin{equation}\label{eq:A3uniformquadratic}
 \#\{u\text{ primitive}:Q(u)=h,\ \hgt(u)\le X\}
 \ll_\eta |h|^2(AX)^\eta\qquad(\eta>0).
\end{equation}
To track $h$, split into the $4h^2$ residue classes modulo $2h$.
The injective map in that proof sends each nonempty class into
$p^2-\operatorname{disc}(Q)r^2=1$, $|p|,|r|\le\lceil6AX^2\rceil$.
This has $O(1+\log(AX))$ solutions with absolute constant: negative or
square discriminants give bounded counts, and positive nonsquare Pell
solutions grow geometrically with ratio at least $2$ in their positive
rational coordinates. Absorb the logarithm to obtain \eqref{eq:A3uniformquadratic}.

For nonzero fixed $x$, the cubic in $y$ is
\[
 16\{\sigma_+(\iota_0x)^2\sigma_+(\iota_0y)^3+
       \sigma_-(\iota_0x)^2\sigma_-(\iota_0y)^3\}.
\]
Its two coefficients are nonzero and its two linear forms independent,
so it has three distinct projective roots. Divide by its positive
content $c$; a nonempty fibre requires $c\mid h$. Hypothesis $\ST$
then bounds primitive solutions by $2800\,3^{1+\omega(|h/c|)}
\ll_\eta|h|^\eta$, uniformly in the cubic, using
$3^{\omega(m)}\le\tau(m)^2$ and the divisor bound (absorbing a sign
into the form when necessary).

For fixed primitive $y$, expansion gives quadratic coefficient height
$\ll(1+|d|)^2\hgt(y)^3$ and discriminant
\[
 \begin{cases}
 -16d((2y_0+y_1)^2-dy_1^2)^3,&d\equiv1\pmod4,\\
 -4096d(y_0^2-dy_1^2)^3,&d\not\equiv1\pmod4.
 \end{cases}
\]
It is nonzero because a nonzero integral vector has nonzero norm.
Since $|d|\le5L^2$ and all coordinates are at most $4Z$, the heights
and ranges are bounded by fixed powers of $P$. Choosing $\eta$
sufficiently small in \eqref{eq:A3uniformquadratic} bounds this fibre
by $O_\delta(h^2P^\delta)$; the cubic fibre is $O_\delta(h^\delta)$.
Both are $O_\delta(h^{2+\delta}P^\delta)$. Summing in either order,
with at most $9D$ choices of $y$ or $9E$ choices of $x$, gives
\[
 \#\{(y,x)\}\ll_\delta h^{2+\delta}P^\delta\min(D,E)
 \ll_{\delta,C_0}L^{56+29\delta}Z^\delta\min(D,E).
\]
Thus $B_{\mathrm f}=56+29\delta$ suffices.
\end{proof}

\begin{proposition}[Graph cover]\label{prop:A3graph}
For every integer $\varkappa_0\ge0$ there are $C_{\mathrm g}\ge2$ and an
integer $k_{\mathrm g}\ge1$, depending on $\varkappa_0$ and $C_0$, with the
following property. Let $L,Z\ge2$ with
$H_{\mathrm g}:=C_{\mathrm g}L^{k_{\mathrm g}}\le Z$, let $\Delta_0$ be
squarefree and nonsquare with $|\Delta_0|\le5L^2$, and let
$1\le e\le C_0L^{28}$. For each orientation there are at most
$H_{\mathrm g}$ intervals $[a_j,b_j]\subset[-1,1]$ and functions $g_j$ of
class $C^{\varkappa_0+1}$ on
$(a_j-2/H_{\mathrm g},b_j+2/H_{\mathrm g})$, with all derivatives of order
at most $\varkappa_0+1$ bounded by $H_{\mathrm g}$ and
$|g_j'|\ge1/H_{\mathrm g}$ on $[a_j,b_j]$, such that every pair of
$\mathcal P(L,Z,\Delta_0,e)$ in that orientation has slopes with
$\xi\in[a_j,b_j]$ and $|\zeta-g_j(\xi)|\le H_{\mathrm g}/Z$ for some $j$.
\end{proposition}
\begin{proof}
Put $d=\Delta_0$. In a fixed orientation denote the dominant coordinates
by $y_1,x_1$, so $\xi=y_0/y_1$, $\zeta=x_0/x_1$.
Each embedding gives an affine factor $t+\omega$ or $1+\omega t$.
Write these as $A_\pm(s)$ for $y$ and $B_\pm(w)$ for $x$. Then
\[
 \Phi(s,w)=16\{A_+(s)^3B_+(w)^2+A_-(s)^3B_-(w)^2\}
\]
is real, and at a counted point
\begin{equation}\label{eq:A3smalltrace}
 |\Phi(\xi,\zeta)|=\frac e{|y_1|^3|x_1|^2}\le C_0L^{28}/Z.
\end{equation}
Constants below may depend on $C_0,\varkappa_0$, but not on $d,L,Z,e$.

\emph{Nearby zeros and nonvanishing partial derivatives.}
For $d<0$, write $\omega=u+iv$, with $u=0$ or $1/2$,
$v^2\ge1/4$, $v^2\ge u^2$. Both possible affine factors have modulus
at least $1/2$ on $\R$: their squared minima are respectively $v^2$
and $v^2/(u^2+v^2)\ge1/2$.
With $A=A_+$ and $B=B_+=rw+t$, put
$a=\operatorname{Re}(A^3r^2)$, $b=\operatorname{Re}(A^3rt)$,
$c=\operatorname{Re}(A^3t^2)$. Expansion gives
\begin{equation}\label{eq:A3imagdisc}
 \Phi=32(aw^2+2bw+c),\qquad b^2-ac=v^2|A|^6\ge1/256.
\end{equation}
A quadratic $aw^2+2bw+c$ with $b^2-ac=\mathcal D>0$ has a real root
at distance at most $|aw^2+2bw+c|/\sqrt{\mathcal D}$: for $a\ne0$
factor it and use that the farther root is at distance at least half
the root separation $2\sqrt{\mathcal D}/|a|$; for $a=0$ use the
linear equation. Consequently a root $w_0$ satisfies
\begin{equation}\label{eq:A3imagclose}
 |w-w_0|\le|\Phi(s,w)|,
\end{equation}
and $|w_0|\le3/2$ once $Z\ge2C_0L^{28}$.
At a zero, $A^3B^2$ is purely imaginary and
$|\operatorname{Im}(A'/A)|=v/|A|^2$ (likewise for $B$), so
\begin{equation}\label{eq:A3imagpartials}
 \Phi_w^2=4096v^2|A|^6,\qquad
 \Phi_s^2=9216v^2|A|^2|B|^4.
\end{equation}
Thus $|\Phi_w|\ge4$ and $|\Phi_s|\ge6$.

For $d>0$, write $A_\pm=p_\pm+q_\pm s$,
$B_\pm=r_\pm+t_\pm w$. Their coefficients are at most $5L$, and
\[
 |p_+q_--p_-q_+|,\ |r_+t_--r_-t_+|\ge1,
\]
since these are $|\omega_+-\omega_-|=\sqrt d$ or $2\sqrt d$.
By \eqref{eq:arith-11}, two embeddings with ratio bounded by $Q$ each
have modulus at least $\hgt(v)/(2Q)$. Dividing by the dominant
coordinate and using (iii), weakened from $L^4$ to $L^6$, gives
\begin{equation}\label{eq:A3realfactors}
 \eta:=1/(2C_0L^6)\le|A_\pm(\xi)|,|B_\pm(\zeta)|\le U:=5L.
\end{equation}
If $|\Phi(s,w)|<32\eta^5$, the $A$ factors have opposite signs;
otherwise $|\Phi|\ge32\eta^5$. The usual discriminant of
$A_+^3(r_++t_+w)^2+A_-^3(r_-+t_-w)^2$ is therefore
\[
 -4A_+^3A_-^3(r_+t_--r_-t_+)^2\ge4\eta^6.
\]
The same root-distance estimate gives
\begin{equation}\label{eq:A3realclose}
 |w-w_0|\le|\Phi(s,w)|/(16\eta^3).
\end{equation}
A polynomial lower threshold on $Z$ makes this displacement at most
$\min(1/2,\eta/(2U))$. Hence $|w_0|\le3/2$, and all factors there
lie in modulus between $\eta/2$ and $2U$.
At such a zero abbreviate the factors to $A,C,B,D$; using
$A^3B^2=-C^3D^2$ gives
\[
 \Phi_s=48\frac{A^2B^2}{C}(q_+C-q_-A),\qquad
 \Phi_w=32\frac{A^3B}{D}(t_+D-t_-B).
\]
The brackets are the constant determinants above, so both partials
have modulus at least $\eta^4/U$.

\emph{Uniform analytic charts.}
In both cases the polynomial extension of $\Phi$ to complex variables
and its partial derivatives of order at most two are bounded by
$K=CL^5\ge1$ on $|s|,|w|\le3$, by expansion of the five affine factors.
For $d<0$ this means extending its real coefficients, not taking a real
part at complex arguments. Put
\[
 \lambda=\frac1{160C_0^4L^{25}},\qquad
 R_w=\frac\lambda{64K},\qquad R_s=\frac{\lambda^2}{4096K^2}.
\]
Both partials at our zeros are at least $\lambda$. At one such zero write
\[
 \Phi(s_0+z,w_0+v)=a(z)v^2+b(z)v+c(z),\qquad b_0=b(0).
\]
Here $c(0)=0$, $|b_0|\ge\lambda$ and, on $|z|<R_s$,
$|a|\le K$, $|b-b_0|,|c|\le K|z|$ (enlarge $K$ absolutely if needed).
In particular
\[
 |b-b_0|\le|b_0|/16,\quad |4ac|\le|b_0|^2/16,\quad
 |c|\le|b_0|R_w/8,\quad 2R_w|a|\le|b_0|/8.
\]
Thus $\mathcal D=(b^2-4ac)/b_0^2$ satisfies $|\mathcal D-1|\le1/4$.
Take the square root equal to $1$ at zero and set
\begin{equation}\label{eq:A3stablebranch}
 G(z)=-\frac{2c(z)}{b(z)+b_0\sqrt{\mathcal D(z)}}.
\end{equation}
The denominator has modulus at least $|b_0|$, so $G$ is holomorphic,
$G(0)=0$, $|G|\le R_w/4$, and $aG^2+bG+c=0$, even when $a=0$.
It is the unique root with $|v|<R_w$, since
$(v-G)\{a(v+G)+b\}=0$ and the second factor differs from $b_0$
by less than $|b_0|/4$. Uniqueness makes $G$ real on the real axis.
Along $g(s)=w_0+G(s-s_0)$ the second-derivative bounds give
$|\Phi_s|,|\Phi_w|\ge\lambda/2$, whence $|g'|\ge\lambda/(2K)$.
Cauchy's estimates on $|s-s_0|\le R_s/2$ give
\[
 |g(s)|\le2,\qquad |g^{(n)}(s)|\le2n!(2/R_s)^n\quad(n\ge1).
\]
All these bounds and reciprocal radii are polynomial in $L$ for fixed $n$.

\emph{Finite cover.}
Partition $[-1,1]\times[-3/2,3/2]$ into rectangles of widths
at most $R_s/16$ and heights at most $R_w/16$. From each rectangle
containing a relevant zero choose one: every zero if $d<0$, and those
with factors in modulus $[\eta/2,2U]$ if $d>0$.
There are $O((R_sR_w)^{-1})$ choices. At each construct the graph above
and retain $[s_0-R_s/4,s_0+R_s/4]\cap[-1,1]$.
Every relevant zero is within $R_s/16,R_w/16$ of the chosen zero in
its rectangle, so uniqueness places it on that graph. Equations
\eqref{eq:A3imagclose} and \eqref{eq:A3realclose} then put every
counted point within $CL^{46}/Z$ vertically of a selected graph.
Choose $H_{\mathrm g}=C_{\mathrm g}L^{k_{\mathrm g}}$ to dominate the
nearby-zero thresholds, the number of rectangles, $8/R_s$, $2K/\lambda$,
$CL^{46}$ and the derivative bounds through order $\varkappa_0+1$.
This is a finite list of polynomial bounds. The enlarged intervals
then remain in $|s-s_0|<R_s/2$, proving every assertion.
\end{proof}

\begin{proposition}[Good-box count]\label{prop:A3good}
There is an absolute integer $\varkappa_0$
such that for every integer $k\ge1$, every $C_2\ge1$ and every $\delta>0$
there are $C_{\mathrm b},B_{\mathrm b}>0$ and an integer $K_{\mathrm b}$
with the following property. Let $L,Z\ge2$ with
$C_{\mathrm b}L^{K_{\mathrm b}}\le Z$, put $H_2=C_2L^k$ and
$M=\lceil Z^{19/100}\rceil$, let $\Delta_0$ be squarefree and nonsquare with
$|\Delta_0|\le5L^2$, and let $1\le e\le C_2L^{28}$. Let
$[a',b']\subset[-1,1]$ and let $g$ be of class $C^{\varkappa_0+1}$ on
$(a'-2r,b'+2r)$, $r\ge1/H_2$, with derivatives of order at most
$\varkappa_0+1$ bounded by $H_2$ and $|g'|\ge1/H_2$ on $[a',b']$. Let
$D,E>M$ with $D^3E^2\le H_2Z^2$. Then, in each orientation, the number of
pairs of primitive vectors $(y,x)$ with $|y_i|\le\sqrt D$,
$|x_i|\le\sqrt E$, $T_0(y;x)=e$ and slopes satisfying $\xi\in[a',b']$ and
$|\zeta-g(\xi)|\le H_2/Z$ is at most
$C_{\mathrm b}L^{B_{\mathrm b}}Z^{97/400+\delta}$.
\end{proposition}
\begin{proof}
Write $H=H_2=C_2L^k$ and fix
\[
 \mu=19/100,\quad q=4096,\quad J=(3q+1)(2q+1),\quad
 \gamma=19q^3/1440,\quad R=qJ+1.
\]
Since $\rho(\mu)=19/360$, this $q$ satisfies \eqref{eq:qlarge}.
Choose an absolute integer $\varkappa_0$ with
\begin{equation}\label{eq:A3Taylororder}
 \mu(\varkappa_0+1)>qJ+\gamma+1.
\end{equation}
Every threshold below is $Z\ge CH^A$, with $C,A$ fixed before $Z$;
finitely many such thresholds give $Z\ge C_{\mathrm b}L^{K_{\mathrm b}}$
with integral $K_{\mathrm b}$.

\emph{Auxiliary forms and their heights.}
Use centres $s_j=a'+j/M$, $0\le j\le\lfloor M(b'-a')\rfloor$,
and cells $[s_j,\min(s_j+1/M,b')]$. Write
$u=\xi-s_j$, $v=\zeta-g(s_j+u)$, so $|u|\le Z^{-\mu}$,
$|v|\le H/Z$. A polynomial threshold ensures $1/M\le r/2$ and
$H/Z\le1$. Let $T$ be the Taylor polynomial of $g(s_j+u)$ of degree
$\varkappa_0$, whose coefficients have modulus at most $H$.
For $0\le a\le3q$, $0\le b\le2q$,
\[
 P_{ab}(u,v)=(s_j+u)^a(T(u)+v)^b
\]
has coefficients at most $CH^{5q+1}$, and Taylor's theorem and
$|x^b-y^b|\le b\max(|x|,|y|)^{b-1}|x-y|$ give error
\[
 |(s_j+u)^a(g(s_j+u)+v)^b-P_{ab}(u,v)|
 \le CH^{5q+1}|u|^{\varkappa_0+1}.
\]
Discarding terms of total degree greater than $\varkappa_0$ adds at most
$CH^{A_0}Z^{-\mu(\varkappa_0+1)}$: their weights $\mu i+j$ are
at least $\mu(\varkappa_0+1)$, and their number and degrees are fixed.

As in Proposition~\ref{prop:vanish}, expand the normalized evaluation
determinant at any $J$ points in the cell. By Proposition~\ref{prop:surplus},
nonzero coefficient determinants use monomials of total weight greater
than $qJ+\rho(\mu)q^3/2=qJ+2\gamma$. The finitely many ordinary terms
are bounded by $CH^{A_1}Z^{-qJ-2\gamma}$; terms with an error row by
$CH^{A_1}Z^{-\mu(\varkappa_0+1)}$. Thus
\[
 |\det\mathcal E|\le CH^{A_1}Z^{-qJ-\gamma}.
\]
Clearing each row by $y_1^{3q}x_1^{2q}$, of modulus at most
$(D^3E^2)^{q/2}\le H^{q/2}Z^q$, gives an integer determinant at most
$CH^{A_1+qJ/2}Z^{-\gamma}$. A polynomial threshold makes it less than
one, hence zero. The finite integer row set therefore has rank less
than $J$ and entries at most $H^{q/2}Z^q$. The signed-minor construction
of Proposition~\ref{prop:aux}, including rank zero, gives a nonzero
integer null vector of height at most $J!H^{qJ/2}Z^{qJ}\le Z^R$ after
another polynomial threshold. Homogenise to a form $G_j$ of bidegree
$(3q,2q)$ vanishing on the cell, with height at most $Z^R$; empty cells
need no form.

\emph{Uniform component estimates.}
Suppose $P\ge2$, $|d|\le5P^2$, and two unimodular substitutions,
an auxiliary form of degree at most $5q$, and the two larger box lengths
all have entries or heights at most $P^v$, for fixed $v$.
We claim the uniform version of Proposition~\ref{prop:components}:
\begin{equation}\label{eq:A3uniformprice}
 C_{v,\delta}e^{2+\delta}P^\delta
 \{1+(\Lambda_1^3\Upsilon_1^2)^{1/10}\},\qquad e\ge1.
\end{equation}
Factor the auxiliary form over $\Q$. There are at most $5q$ factors,
whose heights are $O(P^v)$ by $\GE$. Pure factors of degree at least
two have no primitive zeros; a pure linear factor fixes its primitive
pair up to sign. The uniform quadratic and cubic fibres proved in
Proposition~\ref{prop:A3fallback} then give $O(e^{2+\delta}P^\delta)$:
all transformed heights and ranges are polynomial in $P$, and their
small exponents are chosen after $v,\delta$.

For a mixed factor of bidegree $(\alpha,\beta)$ with
$\min(\alpha,\beta)=1$, the resultant and squarefree-kernel argument
of Proposition~\ref{prop:low} gives
\[
 \#\text{solutions}\le11200d_R\tau(|\Xi|)^{d_R+1},\quad
 |\Xi|\le C|e|\,\|G\|^{12\beta+8\alpha},\quad d_R=3\beta+2\alpha.
\]
Since the possible bidegrees are fixed, the divisor bound gives
$O(|e|^\delta P^\delta)$ with uniform constant; this includes
$(4,1),(1,3)$ and all other factors of minimum bidegree one.
The remaining factors have $\alpha,\beta\ge2$, hence
$10\le2\alpha+3\beta\le3\deg G$. Lemma~\ref{lem:P} gives
$O(P^\delta\{1+(\Lambda_1^3\Upsilon_1^2)^{1/10}\})$, using
$\log\max(1,\Lambda_1^3\Upsilon_1^2)\le5v\log P$.
Its constant is uniform in the transformed trace form and nonzero
right-hand side. Adding the factors proves \eqref{eq:A3uniformprice}.

Taking $P=C_*LZ$, $C_*\ge\max(1,C_2)$, and $e\le C_2L^{28}$ gives
\begin{equation}\label{eq:A3cellprice}
 CL^{56+29\delta}Z^\delta
 \{1+(\Lambda_1^3\Upsilon_1^2)^{1/10}\}.
\end{equation}
Here $v$ changes the constant, not the displayed powers, because all
fibre, divisor and logarithm exponents were chosen after $v,\delta$.

\emph{Reduced bases.}
Impose $M\ge3H$, for instance by $Z\ge(3H)^6$ since $6\mu>1$.
On a nondegenerate interval $g'$ has constant sign, so the centres
$s_j$ are $1/M$-separated and $t_j=g(s_j)$ are $1/(HM)$-separated.
There are at most $5M$ occupied cells, with $|s_j|\le1$, $|t_j|\le2$,
since $|\zeta-t_j|\le H/Z+H/M\le1$. A degenerate interval has a
single cell and the same estimates, with separation vacuous.
Put
\[
 Q=\lceil H\rceil+1\le3H,\quad W=3H,\quad
 M'=QM,\quad D'=W^2Q^2D,\quad E'=W^2Q^2E.
\]
Both sets of centres are $1/M'$-separated. The cell inequalities
\[
 |y_1|\le\sqrt D,\quad|y_0-s_jy_1|\le\sqrt D/M,\qquad
 |x_1|\le\sqrt E,\quad|x_0-t_jx_1|\le W\sqrt E/M
\]
place the points in strips of parameters $(D',M')$, $(E',M')$, because
$\sqrt{D'}/M'=W\sqrt D/M$, and similarly for $E$.
Their reduced integral bases give coordinate bounds
$\Lambda_{1j}\ge\Lambda_{2j}>0$, $\Upsilon_{1j}\ge\Upsilon_{2j}>0$.
The estimates of Section~\ref{sec:sum} give
\[
 \|B_{yj}\|,\|B_{yj}^{-1}\|\le CW^2Q^2D,\qquad
 \Lambda_{1j}\le10WQ\sqrt D,
\]
and analogously for $x,E,\Upsilon$. Since $D,E>M\ge1$ and
$D^3E^2\le HZ^2$, we have $D,E\le HZ^2$. Thus the matrix entries
are $O(H^5Z^2)$ and lengths $O(H^3Z^2)$. Substituting into $G_j$
preserves its bidegree and bounds its height by
$C_qZ^R(CH^5Z^2)^{5q}$. All data in \eqref{eq:A3cellprice} are therefore
bounded by a fixed power of $C_*LZ$. Unimodularity preserves primitivity
and injectivity, so that estimate applies in every cell.

\emph{Cell sum.}
The centres are in a fixed bounded interval, are $1/M'$-separated,
number at most $5M'$, and $D',E'>M'$. Lemma~\ref{lem:moments} gives
\begin{align*}
 \sum_j\Lambda_{1j}^{3/10}\Upsilon_{1j}^{1/5}
 &\le CM'(D'/M')^{3/20}(E'/M')^{1/10}\\
 &=CW^{1/2}Q^{5/4}M(D/M)^{3/20}(E/M)^{1/10}\\
 &\le CW^2Q^2(D^3E^2)^{1/20}M^{3/4}
 \le81CH^{4+1/20}Z^{1/10}M^{3/4}.
\end{align*}
Summing \eqref{eq:A3cellprice} therefore gives
\[
 \#\{(y,x)\}\le CL^{56+29\delta}Z^\delta H^5
 (M+Z^{1/10}M^{3/4})\le4CL^{56+29\delta}H^5Z^{97/400+\delta},
\]
since $M\le2Z^{19/100}$ and $19/100<97/400$.
Thus $B_{\mathrm b}=56+29\delta+5k$ works; enlarge $C_{\mathrm b}$
and choose integral $K_{\mathrm b}$ to enforce all polynomial thresholds.
\end{proof}

\begin{proposition}[Shell count]\label{prop:A3shell}
For every $\eps>0$ there are $C_1,B_1>0$
such that for all $L,Z\ge2$, all squarefree nonsquare $\Delta_0$ with
$|\Delta_0|\le5L^2$ and all $1\le e\le C_0L^{28}$,
\[
 \card\mathcal P(L,Z,\Delta_0,e)\le C_1L^{B_1}Z^{97/400+\eps}.
\]
\end{proposition}
\begin{proof}
Set $\delta=\eps/2$. Choose the fallback constants, then the absolute
Taylor order of Proposition~\ref{prop:A3good}, then its graph-cover
constants $C_{\mathrm g},k_{\mathrm g}$, and finally the good-box
constants with $k=k_{\mathrm g}$, $C_2=\max(128,C_0,C_{\mathrm g})$.
Put $C_{\mathrm{cut}}=C_{\mathrm g}+C_{\mathrm b}+2^{40}$ and
$K_{\mathrm{cut}}=k_{\mathrm g}+K_{\mathrm b}$.
Pairs in $\mathcal P$ have $1\le\hgt(y),\hgt(x)\le2Z$ and
$\hgt(y)^6\hgt(x)^4\le4Z^2$. Lemma~\ref{lem:partition} partitions
this finite set into $O((1+\log Z)^2)$ bins
$D/2<\hgt(y)^2\le D$, $E/2<\hgt(x)^2\le E$, with
$D^3E^2\le128Z^2$ and $\sqrt D,\sqrt E\le4Z$.

If $Z<C_{\mathrm{cut}}L^{K_{\mathrm{cut}}}$, apply the fallback bound
with $D=E=4Z^2$ to obtain
\[
 \card\mathcal P\le4K_{\mathrm f}L^{B_{\mathrm f}}Z^{\delta+2}
 \le4K_{\mathrm f}C_{\mathrm{cut}}^2
 L^{B_{\mathrm f}+2K_{\mathrm{cut}}}Z^\delta.
\]
Otherwise a bin with $\min(D,E)\le Z^{97/400}$ contributes at most
$K_{\mathrm f}L^{B_{\mathrm f}}Z^{97/400+\delta}$.
In every remaining bin,
$M\le2Z^{19/100}\le Z^{97/400}<D,E$, because
$Z\ge2^{40}$ implies $Z^{21/400}\ge2$.
Each of four orientations is covered by at most
$H_{\mathrm g}=C_{\mathrm g}L^{k_{\mathrm g}}$ graphs. The good-box
bound applies with $H_2=C_2L^{k_{\mathrm g}}\ge H_{\mathrm g}$ and
$D^3E^2\le128Z^2\le H_2Z^2$, giving at most
$4C_{\mathrm g}C_{\mathrm b}L^{k_{\mathrm g}+B_{\mathrm b}}
 Z^{97/400+\delta}$ pairs per bin.
Absorb $(1+\log Z)^2\ll_\eps Z^{\eps/2}$ and take
$B_1=2B_{\mathrm f}+3k_{\mathrm g}+B_{\mathrm b}+2K_{\mathrm b}$.
\end{proof}

\subsection{Proof of Theorem~\ref{thm:uniform}}\label{sec:A3proof}
Use Proposition~\ref{prop:A3shell} with $\eps=1/1600$,
so its exponent is $\alpha_0=97/400+1/1600=389/1600$, and set
$\nu=B_1+117$. For $X\ge4C_0L^{10}$, Proposition~\ref{prop:A3lift}
assigns each counted $n\in(X,2X]$ a triple $(e,y,x)$ injectively.
There are at most $C_0L^{28}$ choices of $e$. With
$Z_1=X/(C_0L^9)\ge2$, the height product lies in
$[Z_1,2C_0L^{40}X]$, hence in $[jZ_1,2jZ_1]$ for some integer
$1\le j\le(4C_0^2+2)L^{49}$.
Since $(jZ_1)^{\alpha_0}\le(4C_0^2+2)L^{40}X^{\alpha_0}$, the shell
bound gives
\[
 \#\{X<n\le2X:f(n)\text{ powerful}\}
 \le C_0C_1(4C_0^2+2)^2L^{B_1+28+49+40}X^{\alpha_0}
 =K_{\mathrm w}L^\nu X^{\alpha_0}.
\]
Lemma~\ref{lem:prefix}, with $X_0=4C_0L^{10}$, $n_1=1$ and
$C_{\mathrm w}=K_{\mathrm w}L^\nu$, yields for $N\ge1$
\[
 S(a,b,c;N)\le\left(8C_0L^{10}+
 \frac{K_{\mathrm w}L^\nu}{2^{\alpha_0}-1}+1\right)N^{\alpha_0}
 \le KL^\nu N^{389/1600},\quad
 K=8C_0+\frac{K_{\mathrm w}}{2^{\alpha_0}-1}+1,
\]
because $\nu\ge10$.

For the coefficient sum take $\kappa_1=1/(1600(\nu+3))$,
$M_0=\lfloor N^{\kappa_1}\rfloor$, $N\ge2$, and apply this bound
with $L=2N^{\kappa_1}$. Each term is at most
$K2^\nu N^{\kappa_1\nu+389/1600}$ and there are at most
$M_0(2M_0+1)^2\le9N^{3\kappa_1}$ triples. Since
$\kappa_1(\nu+3)=1/1600$, the sum is at most
$9K2^\nu N^{390/1600}$. \qed

\subsection{The exponent and the constants}\label{sec:A3remarks}
The identity
\[
 \frac{389}{1600}=\frac1{10}+\frac34\cdot\frac{19}{100}+\frac1{1600}
\]
separates the cell-sum exponent from the fixed slack absorbing divisor
bounds and bins. The first two terms exceed $293/1220$ by $57/24400$.
Choosing $\mu_0=19/100\in(57/305,1/5)$ fixes absolute determinant
and Taylor orders and permits polynomial dependence on $L$.
The exponent is convenient, not optimal. No useful numerical value of
$\nu$ or the resulting averaging range is claimed; the point is a fixed
$\eta>0$ permitting a family of coefficients growing with $N$.
\section{Sums of two squares with a varying shift: proof of Theorem~\ref{thm:family}}\label{sec:A2}
We adapt Sections~\ref{sec:charts}--\ref{sec:sum} to $\Z[i]$, choosing
every constant before $a$. Throughout $a\ge1$; the result depends only
on $|a|$. Theorem~\ref{thm:uniform} is not used.

\subsection{The equation}
\begin{lemma}[Gaussian labels]\label{lem:A2labels}
Let $n\ge1$ with $n^2+a^2$ powerful. There are an integer $\ell\ge1$ dividing
$a$, a Gaussian integer $\gamma$ with $1\le\Nm\gamma\le a^2$, and primitive
$y,x\in\Z^2$ such that
\begin{equation}\label{eq:A2}
 n+ia=\ell\,\gamma\,(y_0+iy_1)^3(x_0+ix_1)^2 .
\end{equation}
For given $a$ the pair $(\ell,\gamma)$ takes at most $4\tau(a)^3$ values, and
$(\ell,\gamma,y,x)$ determines $n$.
\end{lemma}
\begin{proof}
Let $P$ divide $n+ia$ exactly once in $\Z[i]$, above $p$.
If $p$ is inert, then $p\mid a$; if $p=2$, the norm has valuation
one, contradicting powerfulness. If $p$ splits, powerfulness forces
the conjugate prime also to divide $n+ia$, so again $p\mid a$.
Take $\gamma$ to be a unit times the product of these $P$.
Its norm divides $\operatorname{rad}(a)^2$, and there are at most
$4\cdot4^{\omega(a)}\le4\tau(a)^2$ choices. Every remaining exponent
$e\ge2$ is $2(e/2)$ if even, and $3+2((e-3)/2)$ if odd; hence
$n+ia=\gamma\mathsf Y^3\mathsf X^2$.
Dividing the coordinates of $\mathsf Y,\mathsf X$ by their gcds
$\ell_1,\ell_2$ gives \eqref{eq:A2}, with
$\ell=\ell_1^3\ell_2^2\mid a$. This adds at most $\tau(a)$ choices.
\end{proof}
No class representatives or balancing are needed. Put
$F_\gamma(y;x)=\operatorname{Im}(\gamma(y_0+iy_1)^3(x_0+ix_1)^2)$.
This integral $(3,2)$-form satisfies $F_\gamma(y;x)=a/\ell$,
while $n$ is $\ell$ times the real part; $\Nm\gamma\le a^2$
controls its coefficients.

\subsection{Small values and one shell}
\begin{lemma}[A quadratic estimate uniform in the right-hand side]
\label{lem:A2quaduniform}
For every $\eta>0$ there is $C_\eta$ such that, for
$H\ge1$, $Z\ge2$, integers $A,B,h$ with
$A^2+B^2\ne0$, $h\ne0$, $|A|,|B|\le H$, the number of primitive
integer pairs $(u,v)$ satisfying
\[
 |u|,|v|\le Z,\qquad B(u^2-v^2)+2Auv=h
\]
is at most $C_\eta(HZ|h|)^\eta$.
\end{lemma}
\begin{proof}
If $B=0$, then $u,v\mid h$, giving at most $4\tau(|h|)^2\ll_\eta|h|^\eta$
pairs. Suppose $B\ne0$. The injective map
$(u,v)\mapsto(x,y)=(Bu+Av,v)$ gives
\[
 x^2-dy^2=Bh,\qquad d=A^2+B^2\ge1,\qquad
 |x|,|y|\le Z_0=\lceil2HZ\rceil\le3HZ.
\]
Vaughan--Wooley~\cite[Lemma~3.5]{vw}, with coefficients $(1,-d,Bh)$,
bounds the solutions with $1\le|x|,|y|\le Z_0$ by
$O_\xi((d|Bh|Z_0)^\xi)$, allowing four sign choices.
The coefficients and right-hand side are nonzero; solutions on the
coordinate axes add at most four. Since
\[
 d|Bh|Z_0\le6H^4Z|h|\le6(HZ|h|)^4,
\]
taking $\xi=\eta/4$ proves the assertion, with constant depending only on $\eta$.
\end{proof}

\begin{proposition}[Fallback]\label{prop:A2fallback}
For every $\eps>0$ there is $K$ such that for all $R_1\ge2$,
all Gaussian integers $\gamma=\gamma'+i\gamma''\ne0$ and all integers
$h\ne0$ with $|\gamma'|,|\gamma''|,|h|\le R_1$, the number of pairs of
primitive vectors with $F_\gamma(y;x)=h$ and $|y|^6|x|^4\le R_1^2$, where
$|y|^2=y_0^2+y_1^2$, is at most $KR_1^{2/5+\eps}$.
\end{proposition}
\begin{proof}
Put $U=|y|^2,V=|x|^2\ge1$. Since $U^3V^2\le R_1^2$,
either $U$ or $V$ is at most $R_1^{2/5}$; each short disc contains
at most $9R_1^{2/5}$ vectors. All coordinates are at most $R_1$.
For fixed $x\ne0$, Lemma~\ref{lem:cubes} makes
$F_\gamma(\cdot\,;x)$ nonzero and squarefree. The content-divisor
argument of Lemma~\ref{lem:cubicfibre} gives at most
$5600\tau(|h|)3^{1+\omega(|h|)}\le16800\tau(|h|)^3
\ll_\eps R_1^\eps$ primitive $y$, uniformly in $x,\gamma,h$.
For fixed $y\ne0$, write $P+iQ=\gamma(y_0+iy_1)^3\ne0$.
The equation in $x$ is $Q(x_0^2-x_1^2)+2Px_0x_1=h$, with
$|P|,|Q|\le8R_1^4$. Lemma~\ref{lem:A2quaduniform}, with
$H=8R_1^4$, $Z=R_1$ and exponent $\eps/6$, gives
$O_\eps(R_1^\eps)$ solutions since $HZ|h|\le8R_1^6$.
Sum cubic fibres over short $x$ and quadratic fibres over short $y$.
\end{proof}

Let $X\ge2$ with $a^2\le X$ and write $a=X^{\upsilon}$, so
$0\le\upsilon\le\frac12$. Define
\begin{equation}\label{eq:m0}
 m_0(\upsilon)=\max\Bigl\{\frac{57+51\upsilon}{305},\
 \frac{47+61\upsilon}{255},\ \frac{4+14\upsilon}{31}\Bigr\}.
\end{equation}
\begin{proposition}[Shell bound]\label{prop:A2shell}
For $0<\delta\le1/100$ and $\eps>0$ there are
$K$ and $X_2$, depending only on $\delta$ and $\eps$, such that for all
$X\ge X_2$, all $1\le a\le\sqrt X$ and all labels $(\ell,\gamma)$, the number
of pairs of primitive $(y,x)$ satisfying \eqref{eq:A2} for some
$n\in(X,2X]$ is at most
\[
 KX^{\eps}\bigl(M+X^{1/10}M^{3/4}\bigr),\qquad
 M=\lceil X^{m_0(\upsilon)+\delta}\rceil .
\]
\end{proposition}
\begin{proof}
\emph{Uniform phase charts.}
For $c^2+d^2=1$ put
\[
 A(s)+iB(s)=(c+id)(s+i)^3,\quad R(s)=(1+s^2)^{3/2},\quad
 P_{c,d}(s,w)=B(s)(w^2-1)+2A(s)w.
\]
Since $A^2+B^2=R^2$, wherever the denominators are nonzero we have
\[
 f_\pm(s)=\frac{\pm B(s)}{R(s)\pm A(s)},\qquad
 P_{c,d}(s,w)=(w-f_\pm(s))(B(s)w+A(s)\pm R(s)).
\]
If $|w|\le1$, $|P_{c,d}(s,w)|\le1/4$ and $D=Bw+A$, then
$D^2=R^2+BP_{c,d}$ and $R\ge1$, $|B|\le R$ imply $|D|\ge R/2$.
Choose the sign of $D$. As $|Bw|\le R$, the corresponding
denominator satisfies $R\pm A\ge R/2\ge1/2$ and
$|D\pm R|\ge1$; hence $|w-f_\pm(s)|\le|P_{c,d}(s,w)|$.

On $[-2,2]$, $|R'|<14$, $|A'|\le21$. Take $h_0=1/1000$ and
centers $t_j=-1+j/1000$, $0\le j\le2000$. Retain \emph{each}
root formula whose denominator at $t_j$ is at least $1/4$, on
$[t_j-h_0,t_j+h_0]$. These at most $4002$ charts are smooth on
$(t_j-3h_0,t_j+3h_0)$: there the denominator exceeds
$1/4-105/1000>1/8$, and $|f_\pm|<96$ since $|B|\le R<12$.
Every $s\in[-1,1]$ is within $h_0/2$ of a center, where the
residual-selected denominator is still at least $1/4$.
Thus the charts cover the stated small-residual set uniformly.

At a root $(Bw+A)^2=R^2>0$, so $\partial_wP\ne0$ and
implicit differentiation gives
\begin{equation}\label{eq:A2phaseode}
 f'(s)=-\frac32\frac{1+f(s)^2}{1+s^2}.
\end{equation}
In particular $|f'|\ge3/10>1/4$ on the enlarged intervals.
Define $H_0(s,w)=w$ and
$H_{k+1}=\partial_sH_k-\frac32(1+w^2)(1+s^2)^{-1}\partial_wH_k$.
The chain rule gives $f^{(k)}(s)=H_k(s,f(s))$; maxima of these
smooth functions on $[-2,2]\times[-100,100]$ give phase-independent
bounds $M_\varkappa$ for $k\le\varkappa+1$.

\emph{The arithmetic tube.}
Rotate $y,x$ by units so that their dominant coordinates
$d_y=y_1,d_x=x_1$ are positive, absorbing inverse units into
$\gamma$. There are at most $16$ cases, preserving $\Nm\gamma$.
With $G=|\gamma|$ and $c+id=\gamma/G$, the slopes $s,w\in[-1,1]$
satisfy
\begin{equation}\label{eq:A2normalizedtube}
 P_{c,d}(s,w)=\frac{a}{\ell Gd_y^3d_x^2}.
\end{equation}
Indeed, take imaginary parts in \eqref{eq:A2}. Dominance gives
$n\le\ell G|y|^3|x|^2\le\sqrt{32}\,\ell Gd_y^3d_x^2$,
so $|P_{c,d}(s,w)|\le8a/X$. For uniformly large $X$, the charts
therefore give $|w-f(s)|\le8X^{-(1-\upsilon)}$, since $a\le\sqrt X$.

\emph{Uniform determinant orders.}
Put $\sigma=1-\upsilon$, $\mu=m_0(\upsilon)+\delta$.
Then $\mu\ge57/305+\delta$, $\sigma\ge1/2$, $\mu<\sigma$ and
$t_0=\mu/\sigma\in(57/305,1)$. Multiplying the three bounds of
Proposition~\ref{prop:A2weights} by $\sigma$ gives leading coefficients
\[
 (305\mu+51\sigma)/18,\quad(255\mu+61\sigma)/18,\quad
 (31\mu+14\sigma)/3.
\]
Each is at least $6+(31/3)\delta$ by \eqref{eq:m0}.
Choose $q\ge1$, $q>(c_5+6)/\delta$, and then $\varkappa$ with
$(57/305+\delta)(\varkappa+1)>n_qq+\delta q^3+1$.
Thus every nonsingular minor has two-scale weight at least
$(6+(31/3)\delta)q^3-c_5q^2>n_qq+\delta q^3$.
Both orders depend only on $\delta$.
On a cell, $|u|\le X^{-\mu}$, $|v|\le8X^{-\sigma}$.
Lemma~\ref{lem:taylor} gives the expansion of
Proposition~\ref{prop:vanish} with weight $\mu i+\sigma j$;
discarded terms have weight at least $\mu(\varkappa+1)$ because
$\mu<\sigma$. The main minors have surplus $\delta q^3$ and
the error a larger saving. Clearing denominators costs at most
$(16X)^q$ per row, since $|y|^6|x|^4\le5X^2$.
Every integer determinant therefore vanishes beyond a uniform
threshold. Proposition~\ref{prop:aux}'s cofactor construction gives
a nonzero $(3q,2q)$-form of height at most $X^{qn_q+1}$.

\emph{All scales and component bounds.}
There are $O((1+\log X)^2)$ dyadic norm boxes, with $D^3E^2\ll X^2$.
The derivative bounds and $X^{-\sigma}\le X^{-\mu}\le2/M$
give strips with one absolute constant. For any $D>0$, apply
Proposition~\ref{prop:bridge} at $R_D=D+M+1>M$.
The norms differ by a scalar, so the same basis is reduced and
$\Lambda_i(D)=\sqrt{D/R_D}\Lambda_i(R_D)$.
Its entries and inverse entries are $O(D+M+1)$ and
$\Lambda_1(D)\le2K_0\sqrt D$; likewise for $E$.
Thus all entries, coordinate lengths, and the transported auxiliary
form and its factors have height at most $X^P$, for fixed $P=P(\delta)$,
using substitution and $\GE$ and increasing the uniform threshold.

Here is the required uniformity in the component count. Put $R=X^P$,
also bounding $|h|=a/\ell$ and the coefficients of $\gamma$.
Transported original coordinates are at most $Z=2R^2$; the parts
of $\gamma(y_0+iy_1)^3$ are at most $H=8RZ^3$, so
$HZ|h|\le128R^{10}$. Lemma~\ref{lem:A2quaduniform} with exponent
$\eta/(10P+1)$ gives $O_{\delta,\eta}(X^\eta)$ per quadratic fibre;
the cubic fibres give the same bound by $16800\tau(|h|)^3$.
Pure lines fix one primitive vector up to sign; higher-degree
irreducible pure factors have no primitive rational zero.
For the four low mixed bidegrees use \eqref{eq:lowuniformrhs}:
the complex presentation survives unimodular substitution, and
$|h|$ and factor height are polynomially bounded in $X$.
For high mixed factors, Lemma~\ref{lem:P} is uniform in the
transported form and $h\ne0$. Summing at most $5q$ factors gives
the per-cell bound
$O_{\delta,\eta}(X^\eta(1+\Lambda_1^{3/10}\Upsilon_1^{1/5}))$.

\emph{Summation.}
Use the regular $1/M$ grid on each branch. Its graph values are
$1/(4M)$-separated by \eqref{eq:A2phaseode}.
The mixed moment extends to all $D,E>0$: apply it to the same
reduced bases with both scales $R=M+1$, then multiply by
$(D/R)^{3/20}(E/R)^{1/10}$. The powers of $R$ cancel, giving
\[
 \sum_{\mathcal C}\Lambda_1(\mathcal C)^{3/10}\Upsilon_1(\mathcal C)^{1/5}
 \ll M(D/M)^{3/20}(E/M)^{1/10}\ll X^{1/10}M^{3/4}.
\]
The constant term contributes $O(M)$. Sum over the bounded number
of branches and rotations, and over the dyadic boxes; choose
$\eta=\eps/2$ and absorb the logarithms in $X^{\eps/2}$.
All constants and thresholds depend only on $\delta,\eps$.
\end{proof}

\subsection{The weight of a minor with two scales}
Give $u^iv^j$ weight $t_0i+j$, $0<t_0<1$. For $g\in\R[u]$ let
$w_q(t_0)$ be the least total weight of monomials whose coefficient
matrix in the shifts $\mathsf b_j^g$ is nonsingular, where
$(\mathsf b_j)$ is a basis of the $(3q,2q)$ rectangle, as in
Proposition~\ref{prop:surplus}.

\begin{proposition}\label{prop:A2weights}
There is an absolute $c_5$ such that for every $g\in\R[u]$ with
$g'(0)\ne0$ and every integer $q\ge1$:
\begin{itemize}
\item[(i)] if $\frac{57}{305}<t_0\le\frac15$, then
$w_q(t_0)\ge\frac{305t_0+51}{18}q^3-c_5q^2$;
\item[(ii)] if $\frac13\le t_0<1$, then
$w_q(t_0)\ge\frac{31t_0+14}{3}q^3$;
\item[(iii)] if $\frac15\le t_0\le\frac13$, then
$w_q(t_0)\ge\frac{255t_0+61}{18}q^3-c_5q^2$.
\end{itemize}
\end{proposition}
\begin{proof}
(i) The computation in Proposition~\ref{prop:surplus} gives
$(17/6+305t_0/18)q^3-O(q^2)$ uniformly on
$57/305<t_0<1/5$. The weight of each fixed family is affine in
$t_0$, so its bound passes to $t_0=1/5$.

(ii) Write $g=g_0+g_1u+g_2u^2+\cdots$, $g_1\ne0$, and put
\[
 W_1=v-g_0,\quad G_1=(1-(g_2/g_1)u)(v-g_0)-g_1u,\quad
 P_{a,b}=u^aW_1^{b-j}G_1^j,\quad j=\min(3q-a,b),
\]
for $0\le a\le3q$, $0\le b\le2q$. Shifting by $g$ gives
$W_1^g=v+g_1u+O(u^2)$ and
$G_1^g=v(1-(g_2/g_1)u)+O(u^3)$, of weighted orders $t_0$ and
at least $1$ because $3t_0\ge1$. Thus $P_{a,b}$ lies in the
rectangle and its shift has order at least $t_0(a+b-j)+j$.

These polynomials form a basis. Use the invertible shift by the tangent
$g_0+g_1u$ and an irrational auxiliary weight $\lambda i+j$,
$1/2<\lambda<1$. Then $W_1\mapsto g_1u+v$ and
$G_1\mapsto v-g_2u^2-(g_2/g_1)uv$, so $P_{a,b}$ has unique initial
term $g_1^{b-j}u^{a+b-j}v^j$. Its exponent $(i,j)$ recovers
$a=\min(3q,i+j)-j$, $b=i+j-a$; the initial weights are therefore
distinct. The least-weight term in any nonzero linear relation
cannot cancel. Independence of these $n_q$ polynomials proves the claim.
Their coefficient matrix at any admissible monomial family is
nonsingular. A nonzero determinant term matches each polynomial
to a monomial of at least its order, giving
\[
 \sum_iw(m_i)\ge\sum_{a,b}\bigl(t_0(a+b-j)+j\bigr)
 =\frac{31t_0+14}{3}q^3+\Bigl(\frac{19t_0}{2}+3\Bigr)q^2
 +\Bigl(\frac{13t_0}{6}+\frac13\Bigr)q.
\]

(iii) Write $t_0=c'/5+c''/3$, $c',c''\ge0$, $c'+c''=1$.
For each admissible family, affinity of its weight and (i),(ii) give
at least $c'(56q^3/9-c_5q^2)+c''73q^3/9$.
Since $(56c'+73c'')/9=(255t_0+61)/18$, the claimed bound follows.
\end{proof}
The chart slope bound supplies $g'(0)\ne0$. The three leading
coefficients after multiplication by $1-\upsilon$ exceed $6$
precisely at the three thresholds in \eqref{eq:m0}; the middle
range is the interpolation in (iii).

\subsection{The exponent and the assembly}
Put $\Theta_0(\upsilon)=1/10+3m_0(\upsilon)/4$ for
$0\le\upsilon\le1/2$. Since
$m_0(\upsilon)+\delta\le11/31+1/100<2/5$, the ceiling bound
$M\le2X^{m_0(\upsilon)+\delta}$ gives
$M\ll X^{1/10}M^{3/4}$. Proposition~\ref{prop:A2shell} is therefore
$\ll X^{\Theta_0(\upsilon)+3\delta/4+\eps}$ per label and shell.
The three terms defining $\Theta_0$ are
\[
 \frac{293+153\upsilon}{1220},\qquad\frac{81+61\upsilon}{340},\qquad
 \frac{61+105\upsilon}{310}.
\]
Their maximum switches at $1/28$ and $19/73$, where
$m_0(\upsilon)/(1-\upsilon)=1/5$ and $1/3$, respectively.
Hence $\Theta_0=\Theta$ on $[0,1/2]$.

\begin{lemma}[Perspective inequality]\label{lem:perspective}
If $0\le\beta_1\le\kappa\le\frac12$ and $0<t_1\le1$ with $\beta_1/t_1\le\frac12$,
then $t_1\,\Theta_0(\beta_1/t_1)\le\Theta_0(\kappa)$.
\end{lemma}
\begin{proof}
Each term of $\Theta_0$ is $c+c'\upsilon$ with $c,c'\ge0$, and
$t_1(c+c'\beta_1/t_1)=ct_1+c'\beta_1\le c+c'\kappa$.
Take their maximum.
\end{proof}

\begin{proof}[Proof of Theorem~\ref{thm:family}]
Choose $\delta$ and the auxiliary $\eps$ small enough for the desired
final exponent. Fix $1\le a\le N^\kappa$ and a label.
For $n\le\max(2a^2,2X_2)$, Proposition~\ref{prop:A2fallback} applies
with $h=a/\ell$ and $R_1\asymp a^2+X_2$: this dominates
$|\gamma|$, $|h|$ and $|y|^3|x|^2\le\sqrt{n^2+a^2}$.
It gives $\ll N^{4\kappa/5+\eps}$, and
$4\kappa/5\le\Theta(\kappa)$ in the stated range.
The remaining $n\le N$ lie in $O(1+\log N)$ shells with $a^2\le X$.
For each use Proposition~\ref{prop:A2shell} and
Lemma~\ref{lem:perspective} with
$t_1=\log X/\log N$, $\beta_1=\log a/\log N$, obtaining
$X^{\Theta_0(\upsilon)}\le N^{\Theta(\kappa)}$.
Lower shells may have $\upsilon>\kappa$, which is why all three
terms of \eqref{eq:m0} are needed. Sum the shells and
$4\tau(a)^3\ll N^\eps$ labels, using
Lemma~\ref{lem:A2labels} to recover $n$.
\end{proof}
The endpoint $289/1050$ is where $\Theta=29/100$; the proof uses
the restriction on $\kappa$ only for
$4\kappa/5\le\Theta(\kappa)$ and $\kappa\le1/2$.

\section{Cubic applications}\label{sec:A4}
\begin{proof}[Proof of the cube-coefficient bound in Corollary~\ref{cor:cubic}]
Put $u=at$, $Q(u)=u^2-u+1$. Since
\[
 u^3+1=(u+1)Q(u),\qquad Q(u)=(u-2)(u+1)+3,
\]
we have $\gcd(u+1,Q(u))\mid3$. If $u^3+1$ is powerful, every prime
$p\ne3$ dividing $Q(u)$ occurs there to order at least two; multiplying
by $9$ also ensures this at $3$. Thus $9Q(u)$ is powerful.
This fixed positive quadratic has nonsquare discriminant $-243$.
The injection $t\mapsto at$, followed by dropping $a\mid u$, gives
\[
 \card\{1\le t\le T:1+a^3t^3\textnormal{ powerful}\}
 \le S(9,-9,9;aT)\ll_\eps(aT)^{293/1220+\eps}
\]
by Theorem~\ref{thm:fixed}, with constant independent of $a$.
\end{proof}

\begin{lemma}\label{lem:cubic-classical-fibres}
Put $P_A(t)=1+At^3$, where $A\ge1$ is an integer. For every $\epsilon>0$
there is $C_\epsilon$, independent of $A,q,k$, such that for all positive
integers $q,k$,
\[
 \begin{split}
 \card\{t\in\Z:P_A(t)\in q\Z^2\setminus\{0\}\}
   &\le C_\epsilon A^{178}q^\epsilon,\\
 \card\{t\in\Z:P_A(t)\in k\Z^3\setminus\{0\}\}
   &\le C_\epsilon A^{62}k^\epsilon.
 \end{split}
\]
\end{lemma}
\begin{proof}
Over $\Q$, Evertse--Silverman~\cite[Theorem~1]{es} bounds integral
arguments giving nonzero squares or cubes of a separable polynomial by,
respectively,
\[
 7^{M(4+9s)}4^{\rho_2(L)},\qquad
 17^{M(6+s)}3^{2Ms+\rho_3(L)}.
\]
Here the degree-$M$ field $L$ contains at least three roots in the square
case and two in the cube case; $s$ counts infinity and the primes of the
polynomial discriminant. The class-group ranks satisfy
$4^{\rho_2(L)}\le h_L^2$, $3^{\rho_3(L)}\le h_L$.

Take $L=\Q(\alpha,\zeta_3)$, $\alpha^3=-A^2$, containing the roots
$\alpha\zeta_3^j/A$ of $P_A$. For noncube $A$, the degree-three and
degree-two fields are linearly disjoint, so the order $\Z[\alpha,\zeta_3]$
has absolute discriminant $(27A^4)^2 3^3$. For cube $A$,
$L=\Q(\zeta_3)$. In either case $M\le6$ and $|D_L|\le3^9A^8$.
Minkowski gives an integral ideal of norm at most $c\sqrt{|D_L|}$
in every class, uniformly in degree at most six. Hence
\[
 h_L\le c^2|D_L|\zeta_L(2)\le c^2|D_L|\zeta(2)^6\ll A^8.
\]

The square fibre gives $qP_A(t)=(qv)^2$, with discriminant supported
on $3Aq$, so $s\le2+\omega(A)+\omega(q)$. Similarly the cube fibre
uses $k^2P_A(t)=(kv)^3$ and $s\le2+\omega(A)+\omega(k)$.
Scaling preserves the distinct roots. The two estimates become
\[
 CA^{16}7^{54\omega(A)}7^{54\omega(q)}
 \ll_\epsilon A^{178}q^\epsilon,\qquad
 CA^8(17^6 3^{12})^{\omega(A)+\omega(k)}
 \ll_\epsilon A^{62}k^\epsilon.
\]
Indeed $7^{54\omega(A)}\le A^{162}$ and $17^6 3^{12}<2^{54}$;
for the fibre parameters use $C_1^{\omega(m)}\ll_\epsilon m^\epsilon$
with fixed $C_1$. All constants are uniform in $A,q,k$.
\end{proof}

\begin{proof}[Proof of the general cubic bound in Corollary~\ref{cor:cubic}]
For $1\le t\le T$, $T\ge2$, write a powerful value as
$P_A(t)=e^2d^3\le2AT^3$. With $U=(2AT^3)^{1/5}$, either $e\le U$
or $d\le U$. Apply Lemma~\ref{lem:cubic-classical-fibres} to
$k=e^2$ in the first case and to $P_A(t)=d(ed)^2$, $q=d$, in the second.
Summing gives
\[
 \card\{1\le t\le T:P_A(t)\textnormal{ powerful}\}
 \ll_\xi A^{62}\sum_{e\le U}e^{2\xi}
       +A^{178}\sum_{d\le U}d^\xi
 \ll_\xi A^{178+1/5+2\xi/5}T^{3/5+6\xi/5}.
\]
Take $\xi=\min(5\eps/6,1)$: then $6\xi/5\le\eps$ and
$178+1/5+2\xi/5<180$, proving the first assertion with
$\nu_1=180$ independently of $\eps$.
\end{proof}

\section{Consequences of the uniform bound}\label{sec:consequences}
In this section $K$ and $\nu$ are the constants of Theorem~\ref{thm:uniform}.

\begin{corollary}\label{cor:shift}
For all integers $A,h,H'\ge1$ and all $T\ge2$,
\[
 \card\{1\le t\le T:At^2+h\textnormal{ is powerful}\}\le K(2h)^{\nu}A^{\nu}T^{389/1600},
\]
and, if $AH'$ is not a square, the number of $t\le T$ with $At^2-H'\ge1$
powerful is at most $K(H'+3)^{\nu}A^{\nu}T^{389/1600}$.
\end{corollary}
\begin{proof}
Apply Theorem~\ref{thm:uniform} to $(A,0,h)$ with $L=2Ah$ (for $h=1$, $L=2A$) and to
$(A,0,-H')$ with $L=A+H'+2\le(H'+3)A$. The discriminants are $-4Ah<0$ and
$4AH'$.
\end{proof}
These give the first assertions of Theorems~\ref{thm:positive} and~\ref{thm:reverse}, with
$C_h=\max(1,\nu)$, $K_h=\max(1,K(2h)^{\nu})$,
$\nu_*(H')=\nu$ and $K_*(H')=K(H'+3)^{\nu}$.

\begin{lemma}[Canonical row]\label{lem:row}
Let $k\ge2$. Every $k$-full $n\ge1$ can be written as
$n=\prod_{i=0}^{k-1}a_i^{k+i}$, where for $i\ge1$ the integer $a_i$ is the
product of the primes $p$ with $v_p(n)\equiv i\pmod k$, and the exponent of
$p$ in $a_0$ is $\lfloor v_p(n)/k\rfloor$, reduced by one if $k\nmid
v_p(n)$. The row $(a_0,\dots,a_{k-1})$ determines $n$.
\end{lemma}
\begin{proof}
Let $v_p(n)=ck+j$ with $0\le j<k$. If $j>0$ then $c\ge1$ because $n$ is
$k$-full, and $(c-1)k+(k+j)=ck+j$. If $j=0$ the exponent is $ck$.
\end{proof}

\begin{proposition}\label{prop:uniform-mixed}
Let $s\in\{6,8\}$ and let $h$ be a nonzero integer. Suppose there are
$K',z>0$ such that for all integers $A\ge1$ that are not squares and all
$T\ge2$,
\begin{equation}\label{eq:hyp}
 \card\{1\le t\le T:At^2+h\ge1,\ At^2+h\textnormal{ is powerful}\}\le K'A^{z}T^{389/1600}.
\end{equation}
Then for every $\eps>0$ there is $C$ such that, for every real $Y\ge2$,
the number of $n\le Y$ such that $n$ is $s$-full,
$n$ is not a square, and $n+h$ is powerful is
at most $CY^{1/s-\delta_s(z)+\eps}$. If \eqref{eq:hyp} holds for all $A\ge1$, the
same bound holds without the condition that $n$ is not a square.
\end{proposition}
\begin{proof}
For $Y<4$ the bound is trivial, so let $Y\ge4$. Write
$\delta=\delta_s(z)$, $\kappa_2=1/(1600(z+1))$ and $T=\sqrt Y\ge2$; then
$s(s+1)\delta\le\kappa_2/2$ and $39/320\le1/s-\delta$. Let
$(a_0,\dots,a_{s-1})$ be the canonical row of $n$ and put
\[
 P=\prod_ia_i,\qquad E_1=\prod_ia_i^{i},\qquad H_3=\prod_{i\ge1}a_i^{s+i},
\]
so that $n=P^{s}E_1=a_0^{s}H_3$ and $H_3\le E_1^{s+1}$.

If $P\le Y^{1/s-\delta}$, we count rows: there are at most
$\tau(P)^{s}\ll Y^{\eps}$ rows with a given product $P$, hence
$\ll Y^{1/s-\delta+\eps}$ such $n$. The condition on $n+h$ is not used here.

If $P>Y^{1/s-\delta}$, then $E_1=n/P^{s}\le Y^{s\delta}$ and
$H_3\le Y^{s(s+1)\delta}\le T^{\kappa_2}$. Since $s$ is even,
$n=H_3t^2$ with $t=a_0^{s/2}\le T$, and $H_3$ is not a square if $n$ is not.
For each $H_3$, the number of $t\le T$ with $H_3t^2+h$ powerful is at most
$K'H_3^{z}T^{389/1600}$ by \eqref{eq:hyp}. Summing over the integers
$H_3\le T^{\kappa_2}$ gives
\[
 \le K'T^{\kappa_2(z+1)+389/1600}=K'T^{390/1600}=K'Y^{39/320}
 \le K'Y^{1/s-\delta}. \qedhere
\]
\end{proof}

\begin{proof}[Proof of the second assertions of Theorem~\ref{thm:positive}]
For $h\ge1$, Corollary~\ref{cor:shift} gives \eqref{eq:hyp} for every
$A\ge1$, with $z=\nu$ (respectively $z=C_h$). An $r$-full integer with
$r\ge2$ is powerful, so Proposition~\ref{prop:uniform-mixed} applies.
\end{proof}

\begin{proof}[Proof of the second assertion of Theorem~\ref{thm:reverse}]
The split is according to whether $n$ is a square. If $n$ is not a square,
take $h=-1$: Corollary~\ref{cor:shift} with $H'=1$ (so $L=A+3$) gives
\eqref{eq:hyp} for nonsquare $A$ with $z=\nu$, and
Proposition~\ref{prop:uniform-mixed} applies. Let $n=m^2$. Since $n$ is
$s$-full with $s\ge4$, every prime dividing $m$ divides it at least twice,
so $m$ is powerful. Also $n-1=(m-1)(m+1)$ is powerful and coprime to $m$.
Hence $(m-1)m(m+1)$ is powerful, and $m-2\le\sqrt Y$ is counted by the
width-three theorem of~\cite[Theorem~1.2]{powerful}, which gives
$\ll Y^{77/640+\eps/2}$ such $m$. Finally
$\frac12(\frac{77}{320}+\eps)\le\frac{49}{400}+\eps$ and
$\frac{49}{400}\le\frac1s-\delta_s(\nu)$, by the second term in the
definition of $\delta_s$.
\end{proof}
As used here, the width-three theorem requires $\CC$ and $\GE$ as well as
$\ST$.

\begin{remark}\label{rem:rs}
The proofs bound the number of $n\le Y$ with $n$ $s$-full and $n+h$
powerful for each fixed $h\ge1$, and the number with $n-1$ powerful.
The hypothesis $r<s$ is not used, $r$ enters only through
$r\ge2$, and the constants do not depend on $r$. The restriction to
$s\in\{6,8\}$ enters as follows: $s$ is even, so that $n=H_3t^2$;
$1/s>39/320$; and, for Theorem~\ref{thm:reverse}, $1/s\ge49/400$. Here $39/320$ is half of $390/1600$, $77/640$ is half of $77/320$, and
$49/400$ lies between $77/640$ and $1/8$. In the definition of $\delta_s$ the first term is the
smaller one for every $z>0$, so $\delta_6(\nu)=1/(268800(\nu+1))$ and
$\delta_8(\nu)=1/(460800(\nu+1))$. Since $\nu$ is not explicit, these savings
cannot be evaluated numerically.
\end{remark}

\section{Remarks on the limits of the method}\label{sec:remarks}
\begin{remark}[The exponent $293/1220$]\label{rem:exponent}
For a cell width $X^{-\mu}$ the method gives the exponent
$\frac1{10}+\frac34\mu$, and the sufficient determinant condition proved here is
$\mu>57/305$. So
$293/1220$ is the infimum of the exponents obtained from
Lemma~\ref{lem:schedule} and Proposition~\ref{prop:surplus}. The term
$\frac1{10}$ is the exponent of Lemma~\ref{lem:P} at the least lifted
degree $10$, attained by components of bidegree $(2,2)$. A smaller exponent
for these components, or a better path $(k,l)$ in
Proposition~\ref{prop:surplus}, would lower $293/1220$. The condition
$\mu<1/5$ is used for the contact polynomials.
\end{remark}

\section{Formal verification and authorship}\label{sec:status}
Theorems~\ref{thm:fixed}, \ref{thm:uniform}, \ref{thm:family},
\ref{thm:positive} and~\ref{thm:reverse} have formal proofs in \lean~4
with Mathlib. Their formal statements take the three estimates of
Section~\ref{sec:cited} as hypotheses. These estimates are published
mathematical results; their proofs are outside the formal development.
The written cubic arguments of Section~\ref{sec:A4} are recorded separately
in Appendix~\ref{app:lean}, including the stronger cube-coefficient
consequence of Theorem~\ref{thm:fixed}.

Independent statement comparisons and a second kernel replay check the
formal declarations. The remaining links are mathematical ones: that the
external estimates imply their formal versions, and that the formal
statements express the claims in the paper. These links, and the written
proofs, have been reviewed by AI systems. Formal verification does not
certify novelty or replace review of the exposition. Appendix~\ref{app:lean}
gives the exact source mapping, trust base and reproduction information.

\subsection*{Use of AI systems}
All results, proofs (written and formal), computation scripts and the text
of this manuscript were produced by AI systems. The results and their
proofs come from large language models running as agents in the Jocelyn
research harness, together with \lean; the writing and revision of the
manuscript, its reviews and the resulting corrections were done with AI
coding assistants, operated by the author. The author chose the problem,
operated the harness, set its objectives, decided the content of the
manuscript and commissioned the reviews, and takes responsibility for the
manuscript. All reviews were carried out by AI systems; no human has
reviewed the results.

\appendix
\section{Formal statements and source mapping}\label{app:lean}
The paper uses its own numbering. The identifiers in the shared formal
library are stable names and do not prescribe the order of exposition.
\begin{center}
\begin{tabular}{@{}ll@{}}
\toprule
Paper statement & Formal interface\\\midrule
Theorem~\ref{thm:fixed} & \texttt{A1}\\
Theorem~\ref{thm:uniform} & \texttt{A3}\\
Theorem~\ref{thm:family} & \texttt{A2}\\
Corollary~\ref{cor:cubic}, general bound & first part of \texttt{A4}\\
Theorem~\ref{thm:positive} & \texttt{A12}; \texttt{A11} is its $h=1$ case\\
Theorem~\ref{thm:reverse} & \texttt{A13}\\\bottomrule
\end{tabular}
\end{center}
In \texttt{A3}, the coefficient-average exponent is the particular choice
$\eta=1/(1600(\nu+3))$ used in the proof of
Theorem~\ref{thm:uniform}. The cube-coefficient bound of
Corollary~\ref{cor:cubic} is the written consequence of \texttt{A1} proved
in Section~\ref{sec:A4}; that stronger conclusion has not been separately
replayed. The library also retains the weaker cube-coefficient bound
$C(a^3)^{\nu_2}T^{389/1600}$ in \texttt{A4}. The general cubic estimate
has a formal proof there using Stewart's estimate for non-cube
coefficients and \texttt{A3}, with all three external inputs, for cube
coefficients. The classical proof given here uses Evertse--Silverman and supplies the coarse choice
$\nu_1=180$. This numerical coefficient exponent is not an additional
replayed claim. The written proof of Lemma~\ref{lem:A2quaduniform}
uses Vaughan--Wooley~\cite[Lemma~3.5]{vw}; the formal development
derives the corresponding quadratic-fibre estimate by a direct Pell argument.

\subsection*{Verification data}
{\raggedright
The pins are \lean~4.33.1 and Mathlib commit
\path{0df444a360eaa60ab8c11dca51a86af692955474}.
The formal interfaces lie in
\path{JocelynWorkspace.Formalization.PowerfulCounting.ConsecutiveCounting},
with explicit hypotheses \path{StewartHyp}, \path{CCDNThm3} and
\path{GelfondHyp}. These are theorem arguments, not new axioms.
The standard axiom set is \texttt{propext}, \texttt{Classical.choice}
and \texttt{Quot.sound}.

The independent statement file is
\path{replay/statements/consecutive-counting/Statements.lean}.
The integrated receipt \path{replay/receipts/consecutive-counting.json}
covers 38 statements in the shared repository, including the interfaces
above. The same-source negative control
\path{replay/receipts/consecutive-counting-negative.json} records rejection
of a deliberately false strengthening of \texttt{A13}.
The sources, pinned dependencies, build commands and replay instructions
are available at \url{https://github.com/dannyply/math-frontier}.
\par}

The source comparisons use Stewart~\cite[Theorem~1, p.~795 and its
consequence on p.~796]{stewart}, the affine-curve theorem
of~\cite[Theorem~3]{ccdn}, and the factor-height inequality
of~\cite[Lemma~1.6.11, p.~27]{bg}. Stewart's consequence is used with
$g=|m|$, squarefreeness supplying the nonzero discriminant. The curve
theorem is used for a prime ideal with linear Hilbert function. The
factor-height statement follows by applying the product-height inequality
to the factor and its integral cofactor. These comparisons, like the
manuscript reviews, were performed by AI systems.

The intermediate declaration map is supplied with the repository as
\path{verification/consecutive-counting/manuscript-source-map.json}.
It distinguishes exact statements, composed proof routes and written
refinements, including the more general tail and moment bounds and the
compressed Gaussian label count. Independent AI readers compared the
mathematical statements and formal sources in both directions, fixing
their initial interpretations before comparison. The formal declarations
belong to the dependency closure of the replayed interfaces; this source
review is not a separate kernel replay of each written lemma. All
mathematical proofs needed here are in the manuscript or the cited literature.

\section{A classical coefficient-uniform comparison}\label{app:classical}
We justify the comparison used in Section~\ref{sec:compare}: for the
quadratics in Theorem~\ref{thm:uniform}, the classical estimates give
\begin{equation}\label{eq:classical-quadratic}
 S(a,b,c;N)\ll_\eps L^{39}N^{2/5+\eps}.
\end{equation}
The coefficient power is deliberately coarse. Put $Q(n)=an^2+bn+c$
and $\Delta=b^2-4ac$. Then $0<|\Delta|\le5L^2$, and the quadratic
root field has absolute discriminant at most $20L^2$, hence class
number $\ll L^2$ by the Minkowski argument in
Lemma~\ref{lem:cubic-classical-fibres}.

For $k\ge1$, the equation $Q(n)=kv^3$ gives
$k^2Q(n)=(kv)^3$. The cube estimate of Evertse--Silverman
used in that lemma applies with field degree two and
$s\le1+\omega(\Delta)+\omega(k)$. Since $17^2 3^4<2^{18}$,
its bound is
\[
 \card\{n\in\Z:Q(n)\in k\Z^3\setminus\{0\}\}
 \ll_\xi L^2|\Delta|^{18}k^\xi
 \ll_\xi L^{38}k^\xi.
\]
The roots are distinct because $\Delta\ne0$.

For a square fibre $Q(n)=qv^2$ with $q\ge1$ and $1\le n\le N$,
completion of the square gives
\[
 (2an+b)^2-4aqv^2=\Delta,
 \qquad |2an+b|,|v|\le3LN.
\]
Vaughan--Wooley~\cite[Lemma~3.5]{vw} bounds the solutions by
$C_\xi(L^4qN)^\xi$, uniformly in these parameters. Its nonzero
right-hand-side condition follows from $\Delta\ne0$; signs and a
zero first coordinate change the count only by a constant.

Write each counted powerful value as $Q(n)=e^2d^3\le3LN^2$,
and put $U=(3LN^2)^{1/5}$. Either $e\le U$ or $d\le U$.
Summing the cube fibres with $k=e^2$ gives
$C_\xi L^{38}U^{1+2\xi}$; summing the square fibres with $q=d$
gives $C_\xi L^{4\xi}N^\xi U^{1+\xi}$.
Thus the two contributions are bounded by
\[
 C_\xi L^{38+1/5+2\xi/5}N^{2/5+4\xi/5},\qquad
 C_\xi L^{1/5+21\xi/5}N^{2/5+7\xi/5}.
\]
Taking $\xi=\min(5\eps/7,1/4)$ proves
\eqref{eq:classical-quadratic}. This is a consequence of the cited
estimates, not a claim that those sources state this coefficient power.

\begin{thebibliography}{99}
\bibitem{powerful} D. Plymesser, \emph{Counting powerful products of consecutive integers}, \url{https://doi.org/10.5281/zenodo.23614742}.
\bibitem{stewart} C.~L. Stewart, On the number of solutions of polynomial congruences and Thue equations, \emph{J. Amer. Math. Soc.} \textbf{4} (1991), 793--835.
\bibitem{ccdn} W. Castryck, R. Cluckers, P. Dittmann and K.~H. Nguyen, The dimension growth conjecture, polynomial in the degree and without logarithmic factors, \emph{Algebra Number Theory} \textbf{14} (2020), 2261--2294; \href{https://arxiv.org/abs/1904.13109v1}{arXiv:1904.13109v1}.
\bibitem{bg} E. Bombieri and W. Gubler, \emph{Heights in Diophantine Geometry}, Cambridge University Press, 2006.
\bibitem{hbannals} D.~R. Heath-Brown, The density of rational points on curves and surfaces, \emph{Ann. of Math. (2)} \textbf{155} (2002), 553--595.
\bibitem{hb} D.~R. Heath-Brown, Square-free values of $n^2+1$, \emph{Acta Arith.} \textbf{155} (2012), 1--13.
\bibitem{wm} W. Wongcharoenbhorn and Y. Meemark, Square-full values of quadratic polynomials, \emph{Bull. Aust. Math. Soc.} \textbf{113} (2026), 40--54.
\bibitem{puchta} J.-C. Puchta, On the powerfull part of $n^2+1$, \emph{Arch. Math. (Brno)} \textbf{39} (2003), 187--189.
\bibitem{saito} Y. Saito, \emph{Shared prime support, small radicals, and square-full polynomial values}, unrefereed manuscript, \url{https://doi.org/10.5281/zenodo.21942579}.
\bibitem{es} J.-H. Evertse and J.~H. Silverman, Uniform bounds for the number of solutions to $Y^n=f(X)$, \emph{Math. Proc. Cambridge Philos. Soc.} \textbf{100} (1986), 237--248.
\bibitem{vw} R.~C. Vaughan and T.~D. Wooley, Further improvements in Waring's problem, \emph{Acta Math.} \textbf{174} (1995), 147--240.
\bibitem{ddl} J.-M. De Koninck, N. Doyon and F. Luca, Powerful values of quadratic polynomials, \emph{J. Integer Seq.} \textbf{14} (2011), Article 11.3.3.
\bibitem{pasten} H. Pasten, Powerful values of polynomials and a conjecture of Vojta, \emph{J. Number Theory} \textbf{133} (2013), 2964--2998.
\bibitem{streeter} S. Streeter, Campana points and powerful values of norm forms, \emph{Math. Z.} \textbf{301} (2022), 627--664.
\bibitem{pohl} S.~S. Pohl, $S$-squareful values of binary quadratic forms, \href{https://arxiv.org/abs/2610.20316v1}{arXiv:2610.20316v1}, preprint.
\end{thebibliography}
\end{document}
